\documentclass[11pt]{article}

\usepackage[margin=1in]{geometry}
\usepackage[T1]{fontenc}
\usepackage[utf8]{inputenc}
\usepackage{lmodern}
\usepackage{microtype}
\usepackage{amsmath,amssymb,amsthm}
\usepackage{mathtools}
\usepackage{hyperref}
\usepackage{booktabs}
\usepackage{enumitem}
\usepackage{listings}
\usepackage{xcolor}

\hypersetup{
  colorlinks=true,
  linkcolor=blue!50!black,
  citecolor=blue!50!black,
  urlcolor=blue!50!black
}

\lstdefinelanguage{Lean4}{
  morekeywords={
    import,namespace,end,open,section,noncomputable,abbrev,
    def,theorem,lemma,where,by,fun,let,in,have,show,exact,
    intro,apply,simpa,calc
  },
  sensitive=true,
  morecomment=[l]{--},
  morecomment=[s]{/-}{-/},
  morestring=[b]"
}

\lstset{
  language=Lean4,
  basicstyle=\ttfamily\small,
  keywordstyle=\bfseries\color{blue!50!black},
  commentstyle=\itshape\color{green!40!black},
  stringstyle=\color{red!50!black},
  showstringspaces=false,
  breaklines=true,
  frame=single,
  framerule=0.3pt,
  xleftmargin=0.5em,
  xrightmargin=0.5em,
  keepspaces=true,
  columns=fullflexible,
  literate=
    {¬}{{$\neg$}}1
    {∃}{{$\exists$}}1
    {→}{{$\to$}}1
    {≤}{{$\leq$}}1
}

\newtheorem{theorem}{Theorem}[section]
\newtheorem{lemma}[theorem]{Lemma}
\theoremstyle{remark}
\newtheorem{remark}[theorem]{Remark}

\title{P vs NP Grassroots:\\
Ground-Up Background Theory for a Structured Switching Route}
\author{}
\date{\today}

\begin{document}
\maketitle

\begin{abstract}
This note records the \emph{grassroots} side of the current Lean work around a
proposed route to $P \neq NP$.  It is intentionally not the blocker note.
The companion crux-first analysis lives separately in \texttt{pnp\_crux.tex}.

The goal here is narrower and more constructive: build reusable background theory
from the ground up, using mathlib and existing Mettapedia infrastructure wherever
possible, so that if the high-level proof strategy can be repaired, it lands in a
clean machine-checked framework rather than in ad hoc prose.

The current grassroots state now includes explicit encoded hypothesis classes,
finite empirical-risk minimization, finite deceptive-sample bounds, threshold
recovery lemmas, a rational uniform-rate layer, and explicit uniform success
bounds for exact ERM recovery under finite sampling.  The success layer now
also exposes non-deceptive sample-rate lower bounds directly, so downstream
arguments can use either the realized-target recovery theorem or the
target-independent complement of the deceptive region.  It now also includes
the first weighted finite-sampling steps beyond uniform draws: an exact
one-predictor agreement law and a family-level weighted union bound over finite
encoded classes, together with weighted exact-recovery lower bounds for ERM.
It also now packages the exact \emph{same-route interface}: a repair that still
intends to use switching plus ERM must produce a bit-encoded family over an
available input surface and therefore inherits the current finite weighted
recovery theorem.
It also now isolates the exact post-switch visible surface and proves a clean
finite-summary kernel: any switched family that truly factors through a finite
summary surface automatically inherits an explicit truth-table bit budget equal
to the cardinality of that summary type.
It now also ties the four current exact-surface affine candidate classes
directly to the weighted recovery backbone: once the target lies in one of
those classes, the finite-sampling lower bound is available immediately with the
matching explicit bit budget.
It now also includes a probability-free finite-count conditional-independence
API: exact independence, approximate tolerance, output-fiber certificates, and
finite-output max defects are all connected by small reusable Lean theorems.
It now also turns the balanced-fiber fork into grassroots source-erasure
infrastructure: balanced finite-coin output fibers are exactly neutrality for
all decidable output predicates, and such erasure forces every Boolean output
decoder to have exactly half aggregate accuracy.
This is now connected back to the finite-count independence API: on the sample
space \(\mathrm{Bool}\times\mathrm{Coin}\), predicate neutrality is equivalent
to exact finite-count independence from the source bit, and balanced fibers are
equivalent to zero output defect.
The quantitative form is also explicit: each output fiber's symmetric defect is
exactly \(|\mathrm{Coin}|\) times its true/false coin-fiber imbalance, so an
approximate-independence tolerance is precisely a uniform bound on those scaled
imbalances.
It now also isolates a sharper shared-basis regime: if one fixed affine feature
extractor is reused across the switched family and only the downstream combiner
varies, then the code budget collapses to the combiner alone.
The Meredith-side limit is now considerably sharper.  The currently formalized
weakness bridge exports only global evidence scalars, so any predictor factoring
only through those summaries is constant on the whole state space.  The current
exact HML layer is infinitary by default: if the certified minimal-context type
is infinite, then even the depth-1 fragment is infinite, and the concrete rho
instance already realizes that case.  And one layer below exact HML, the
concrete rho present-moment surface already realizes arbitrary finite
membership patterns on a natural probe slice: a single-output probe state
against a flat bag of input guards records exactly which channels are present in
the environment.
On the semantic side, it now includes a depth-bounded GSLT/HML observational
kernel, a bridge from bisimulation quotients to quantale weakness, and a first
concrete minimal-context instance for the rho calculus.  This isolates the exact
positive shape that a successful switching theorem would need: not mere
locality, but locality together with a concrete compression theorem, a clean
finite learning interface, and an explicit observational/evidence semantics.
\end{abstract}

\tableofcontents

\section{Scope}

The October 2025 manuscript
\emph{``P != NP: A Non-Relativizing Proof via Quantale Weakness and Geometric
Complexity''} proposes a route from a structured Unique-SAT ensemble to a
distributional lower bound, and then to a contradiction with the upper bound that
would follow from $P=NP$.

This document does \emph{not} ask where that route fails.  That is the job of the
companion crux note.  Instead it asks a different question:

\begin{quote}
What background theory would we want in Lean even if the final proof still needs
repair?
\end{quote}

That question leads to a program with three design constraints:
\begin{enumerate}[leftmargin=1.5em]
  \item Prefer general infrastructure over one-off rebuttal lemmas.
  \item Reuse mathlib computability and existing Mettapedia semantics rather than
  inventing parallel foundations.
  \item State optimistic bridge theorems in a way that makes missing hypotheses
  explicit, so later fixes can slot into them cleanly.
\end{enumerate}

\section{Existing Lean Base}

Several ingredients already exist in the local formalization and are useful for a
ground-up PNP effort.

\begin{center}
\begin{tabular}{ll}
\toprule
\textbf{Module} & \textbf{Grassroots role} \\
\midrule
\texttt{Mathlib.Computability.Partrec} & partial recursive computation infrastructure \\
\texttt{KolmogorovComplexity/Basic.lean} & universal algorithm and plain complexity \\
\texttt{Hyperseed/Ultrainfinitism.lean} & observer-relative available-region semantics \\
\texttt{Hyperseed/Basic.lean} & closure/cascade interfaces on accessible observations \\
\texttt{GSLT/Logic/LogicalMetric.lean} & finite-depth observational distinction kernel \\
\texttt{GSLT/Meredith/WeaknessBridge.lean} & operational distinguishability $\to$ weakness/evidence \\
\texttt{GSLT/Meredith/RhoMinimalContext.lean} & first concrete minimal-context / HML anchor \\
\texttt{PNP/HypothesisClass.lean} & encoded classifier families and code-budget bounds \\
\texttt{PNP/FiniteERM.lean} & empirical-risk minimization kernel for finite encoded families \\
\texttt{PNP/FiniteConsistencyBound.lean} & finite uniform consistency bound for deceptive samples \\
\texttt{PNP/FiniteUniformRecovery.lean} & threshold corollaries guaranteeing good samples from cardinality gaps \\
\texttt{PNP/BitFamilyUniformRecovery.lean} & bit-budget exact-recovery wrapper for finite ERM \\
\texttt{PNP/FiniteUniformRecoveryRegression.lean} & regression wrappers for the bit-budget recovery threshold \\
\texttt{PNP/BitFamilyUniformRecoveryBinary.lean} & two-point-domain specialization reducing the threshold to \(s<m\) \\
\texttt{PNP/BitFamilyUniformRecoveryBinaryRegression.lean} & regression wrappers for the binary recovery specialization \\
\texttt{PNP/FiniteUniformRate.lean} & rational uniform-rate bounds for deception and exact recovery \\
\texttt{PNP/FiniteUniformSuccess.lean} & miss-ratio decay, non-deceptive-rate lower bounds, and explicit uniform/bit-budget success bounds \\
\texttt{PNP/FiniteUniformSuccessRegression.lean} & regression wrappers for the bit-budget non-deceptive and recovery bridges \\
\texttt{PNP/BitFamilyUniformSuccessBinary.lean} & Boolean \(1/2\)-miss quantitative non-deceptive/recovery bounds, positive-margin strictness, and binary positivity \\
\texttt{PNP/BitFamilyUniformSuccessBinaryRegression.lean} & regression wrappers for Boolean quantitative non-deceptive and recovery bounds \\
\texttt{PNP/ClockedKpolyBridge.lean} & construction, family realization interface, and finite success-rate/good-sample bridge for clocked \(s\)-bit families \\
\texttt{PNP/ClockedKpolyBridgeRegression.lean} & regression wrappers for the clocked construction, realization, success-rate, and good-sample bridge \\
\texttt{PNP/ClockedKpolyTargetInterface.lean} & repackages exact bit-budget witnesses as clocked families with recovery rates, non-deceptive rates, and good samples \\
\texttt{PNP/ClockedKpolyTargetInterfaceRegression.lean} & regression wrappers for clocked target-interface constructors, rates, and good samples \\
\texttt{PNP/ClockedKpolyBridgeEquivalence.lean} & clocked realization equivalence with exact bit budgets and exact visible compression targets \\
\texttt{PNP/ClockedKpolyBridgeEquivalenceRegression.lean} & regression wrappers for clocked bridge equivalences \\
\texttt{PNP/ClockedKpolyGapAssessment.lean} & bundles the clocked finite-learning payload and proves it is exactly the finite predictor-image obligation \\
\texttt{PNP/ClockedKpolyGapAssessmentRegression.lean} & regression wrappers for the gap-assessment equivalences and obstructions \\
\texttt{PNP/ClockedKpolyActualGapClosure.lean} & closes the clocked payload for the concrete exact ZAB ERM family and records the bare-interface full-family obstruction \\
\texttt{PNP/ClockedKpolyActualGapClosureRegression.lean} & regression wrappers for the actual-family closure and current-interface obstruction \\
\texttt{PNP/ActualSwitchedLocalFamily.lean} & manuscript-shaped local normal form \(output_i(\phi)=h_i(z(\phi),a_i(\phi),b(\phi))\), with ZAB ERM closure and locality-only no-go \\
\texttt{PNP/ActualSwitchedLocalFamilyRegression.lean} & regression wrappers for the actual local-family closure and obstruction \\
\texttt{PNP/ActualSwitchedLocalStructure.lean} & bounded-code and shared \((zfeat(z),a,b)\) decision-list strengthenings of the actual local interface, with finite-image closure and full-rule exclusions \\
\texttt{PNP/ActualSwitchedLocalStructureRegression.lean} & regression wrappers for the strengthened actual local-interface closure and exclusions \\
\texttt{PNP/ActualSwitchedLocalERMConstruction.lean} & concrete actual switched-local ERM wrapper from \(zOf,aOf,bOf,zfeat\), proving ZAB data, bit-code data, finite cover, and clocked payload \\
\texttt{PNP/ActualSwitchedLocalERMConstructionRegression.lean} & regression wrappers for the constructed actual switched-local ERM closure \\
\texttt{PNP/ActualSwitchedLocalCodeConstruction.lean} & concrete actual switched-local wrapper from explicit per-index ZAB decision-list codes, proving the same finite-cover and clocked-payload closure \\
\texttt{PNP/ActualSwitchedLocalCodeConstructionRegression.lean} & regression wrappers for the explicit-code actual switched-local closure \\
\texttt{PNP/ActualSwitchedLocalPluginLookupObstruction.lean} & manuscript plug-in lookup endpoint and full-rule obstruction to treating finite-alphabet lookup as ZAB/code closure \\
\texttt{PNP/ActualSwitchedLocalPluginLookupObstructionRegression.lean} & regression wrappers for the plug-in lookup obstruction \\
\texttt{PNP/ActualSwitchedLocalPluginMajorityObstruction.lean} & counted plug-in majority endpoint and full-rule obstruction via one vote per visible input \\
\texttt{PNP/ActualSwitchedLocalPluginMajorityObstructionRegression.lean} & regression wrappers for the plug-in majority obstruction \\
\texttt{PNP/ActualSwitchedLocalPluginSampleMajorityObstruction.lean} & finite-sample plug-in majority endpoint and full-rule obstruction via the graph sample of each visible rule \\
\texttt{PNP/ActualSwitchedLocalPluginSampleMajorityObstructionRegression.lean} & regression wrappers for the sample-level plug-in majority obstruction \\
\texttt{PNP/ActualSwitchedLocalBoundedSampleMajorityObstruction.lean} & bounded finite-sample plug-in majority endpoint with exact surjectivity threshold at the visible alphabet size \\
\texttt{PNP/ActualSwitchedLocalBoundedSampleMajorityObstructionRegression.lean} & regression wrappers for the bounded-sample majority obstruction \\
\texttt{PNP/ClockedKpolyCompressionObstruction.lean} & surjectivity lower bound for low-bit clocked realization attempts \\
\texttt{PNP/ClockedKpolyCompressionObstructionRegression.lean} & regression wrappers for clocked compression obstructions \\
\texttt{PNP/FiniteIIDAgreement.lean} & exact weighted one-predictor agreement mass under a finite PMF \\
\texttt{PNP/FiniteIIDFamilyBound.lean} & weighted deceptive-mass union bound for finite encoded families \\
\texttt{PNP/FiniteIIDRecovery.lean} & weighted exact ERM recovery lower bounds under finite PMFs \\
\texttt{PNP/SameRouteInterface.lean} & exact certificate for when a repair still counts as the same switching/ERM route \\
\texttt{PNP/BitFamilyWeightedRecovery.lean} & epsilon-form weighted recovery certificate for bit-encoded same-route families \\
\texttt{PNP/BitFamilyWeightedRecoveryRegression.lean} & regression wrapper for the weighted epsilon recovery certificate \\
\texttt{PNP/BadCodeAgreementObstruction.lean} & atomic-measure obstruction to strict bad-code agreement bounds in the exact ZAB recovery route \\
\texttt{PNP/BadCodeAgreementObstructionRegression.lean} & regression checks for the pure-atom bad-code obstruction \\
\texttt{PNP/ExactSwitchedFamily.lean} & exact post-switch surface plus finite-summary compression kernel \\
\texttt{PNP/ExactVisibleImageBudget.lean} & exact and clocked bit budgets as finite predictor-image covers of size at most \(2^s\) \\
\texttt{PNP/ExactVisibleImageBudgetRegression.lean} & regression checks for finite predictor-image budget equivalences \\
\texttt{PNP/ProductSuccessKpolyBridgeObstruction.lean} & product/fielded success facts alone do not imply exact-visible finite-image compression, positive uses reduce to finite-image cover obligations, and already-surjective exact families refute the bridge below threshold \\
\texttt{PNP/ProductSuccessKpolyBridgeObstructionRegression.lean} & regression checks for the product/fielded bridge counterexamples, semantic-boundary equivalences, and surjective-family refuters \\
\texttt{PNP/ABVisibleSurface.lean} & reduced raw visible surface $(a,b)$ and its explicit cardinality \\
\texttt{PNP/AffineColumnFamily.lean} & concrete affine GF(2)-style column family with $k+1$ bit budget \\
\texttt{PNP/AffineFeatureFamily.lean} & bounded affine-feature family with budget $r(k+1)+2^r$ \\
\texttt{PNP/SparseThresholdAffineFamily.lean} & sparse-threshold affine-feature family with linear budget $r(k+3)$ \\
\texttt{PNP/AffineDecisionListFamily.lean} & affine decision-list family with linear budget $r(k+2)+1$ \\
\texttt{PNP/ExactAffineRecovery.lean} & exact-surface weighted recovery bounds for the current affine candidate classes \\
\texttt{PNP/SharedAffineFeatureFamilies.lean} & shared-basis exact-surface families with combiner-only budgets and recovery bounds \\
\texttt{PNP/ExactABDecisionListFamily.lean} & concrete raw exact-surface decision-list family on $(a,b)$ with budget $2k+1$ \\
\texttt{PNP/ABDecisionListRoute.lean} & pullback certificate from reduced $(a,b)$ decision lists to exact-surface compression/recovery \\
\texttt{PNP/ABDecisionListObstruction.lean} & strict-subclass theorem for raw $(a,b)$ fixed-order decision lists versus all raw Boolean rules \\
\texttt{PNP/ABDecisionListObstructionRegression.lean} & regression checks for full raw-rule and exact-surface pullback refuters \\
\texttt{PNP/ExactABAffineFamilies.lean} & affine-feature, sparse-threshold, and decision-list ladders on the raw $(a,b)$ surface \\
\texttt{PNP/SharedExactABFeatureFamilies.lean} & shared-basis affine ladders on the raw $(a,b)$ surface with combiner-only budgets \\
\texttt{PNP/SharedABFeatureRoutes.lean} & quotient-route transfer from reduced $(a,b)$ families to shared-basis exact-surface budgets \\
\texttt{PNP/SharedABFeatureObstruction.lean} & strict-subclass and exact-pullback refuters for shared raw $(a,b)$ feature classes below the reduced truth-table threshold \\
\texttt{PNP/SharedABFeatureObstructionRegression.lean} & regression checks for shared raw $(a,b)$ feature obstruction anchors \\
\texttt{PNP/SharedABTargetInterface.lean} & unified target data packages for the shared raw $(a,b)$ route \\
\texttt{PNP/CanonicalABTargetRoute.lean} & canonical-coordinate specialization reducing the target interface back to raw $(a,b)$ decision lists \\
\texttt{PNP/CanonicalABCandidateInterface.lean} & single candidate-data package for the canonical raw $(a,b)$ decision-list route, with recovery consequence \\
\texttt{PNP/CanonicalABCodeWitness.lean} & explicit code-witness reduction: the route follows automatically from one raw $(a,b)$ decision-list code assignment per index \\
\texttt{PNP/CanonicalABRecoveryInterface.lean} & final generic interface: explicit code assignment plus agreement bound imply compression and recovery \\
\texttt{PNP/BitFamilyERM.lean} & ERM transport layer: one fixed local class yields one selected code per sample/index \\
\texttt{PNP/ExactZABDecisionListFamily.lean} & exact-surface decision-list family on $(zfeat(z),a,b)$ with budget $r+2k+1$ \\
\texttt{PNP/ExactZABTargetInterface.lean} & target and recovery interfaces for the shared $(zfeat(z),a,b)$ decision-list route, including finite-image presentations \\
\texttt{PNP/ExactZABERMRoute.lean} & ERM specialization: the selected code itself witnesses the shared $(zfeat(z),a,b)$ route \\
\texttt{PNP/SharedExactZABFeatureFamilies.lean} & shared-basis affine ladders on $(zfeat(z),a,b)$ with combiner-only budgets \\
\texttt{PNP/SharedExactZABTargetInterface.lean} & unified target and recovery interfaces for the shared-basis $(zfeat(z),a,b)$ route \\
\texttt{PNP/SharedExactZABERMRoute.lean} & ERM specialization of the shared-basis $(zfeat(z),a,b)$ route with combiner-only budget \\
\texttt{PNP/CanonicalZABTargetRoute.lean} & canonical-coordinate specialization reducing the shared-basis exact route back to raw $(zfeat(z),a,b)$ decision lists \\
\texttt{PNP/CanonicalZABCandidateInterface.lean} & concrete candidate package for the direct exact $(zfeat(z),a,b)$ route, including finite-image presentations and recovery consequence \\
\texttt{PNP/CanonicalZABCodeWitness.lean} & explicit code-witness reduction: one exact $(zfeat(z),a,b)$ decision-list code assignment per index suffices and exposes the finite image \\
\texttt{PNP/CanonicalZABRecoveryInterface.lean} & final direct exact-route interface: code assignment plus agreement bound imply finite-image compression and recovery \\
\texttt{PNP/CanonicalZABERMRoute.lean} & ERM-to-canonical bridge: the exact ERM wrapper itself supplies the canonical code witness and recovery interface \\
\texttt{PNP/CanonicalZABERMInterface.lean} & final manuscript-facing ERM interface: samples, equality to the ERM wrapper, and an agreement bound imply compression and recovery \\
\texttt{PNP/BitVecZABVisibleSurface.lean} & when $z$ is already bit-valued, the full exact visible surface $(z,a,b)$ becomes one raw bitvector of length $r+2k$ \\
\texttt{PNP/BitVecZABERMInterface.lean} & concrete identity-extractor specialization of the final exact ERM interface on the raw full visible bit surface \\
\texttt{PNP/ProjectedZABERMInterface.lean} & coordinate-projection specialization: the extractor burden is reduced to choosing finitely many coordinates of a bit-valued $z$ \\
\texttt{PNP/GlobalWeaknessObstruction.lean} & current weakness summaries alone cannot yield a nonconstant predictor surface \\
\texttt{PNP/InfinitaryHMLObstruction.lean} & infinite minimal contexts force an infinite depth-1 HML observable surface \\
\texttt{PNP/PresentMomentShattering.lean} & exact rho present-moment probes already realize arbitrary finite membership slices \\
\texttt{PNP/ConditioningObstruction.lean} & exact finite-count conditional independence and conditioning-collapse examples \\
\texttt{PNP/ApproximateDecorrelationObstruction.lean} & oriented finite-count defects and recoding obstruction applications \\
\texttt{PNP/FiniteCountIndependenceFoundation.lean} & grassroots API for exact, approximate, fiberwise, and finite-output defect certificates \\
\texttt{PNP/FiniteCountIndependenceFoundationRegression.lean} & regression wrappers for the finite-count independence API \\
\texttt{PNP/FiniteCoinSourceErasureFoundation.lean} & balanced finite-coin fibers as predicate erasure and decoder half-accuracy \\
\texttt{PNP/FiniteCoinSourceErasureFoundationRegression.lean} & regression wrappers for finite-coin source-erasure infrastructure \\
\texttt{PNP/FiniteCoinIndependenceBridge.lean} & predicate erasure as exact finite-count source/output independence \\
\texttt{PNP/FiniteCoinIndependenceBridgeRegression.lean} & regression wrappers for the finite-coin independence bridge \\
\texttt{PNP/FiniteCoinDefectFormula.lean} & exact scaled-imbalance formula for finite-coin output defects \\
\texttt{PNP/FiniteCoinDefectFormulaRegression.lean} & regression wrappers for the finite-coin defect formula \\
\texttt{PNP/FiniteCoinToleranceQuantization.lean} & tolerance-below-\(|\mathrm{Coin}|\) collapse to exact balanced fibers \\
\texttt{PNP/FiniteCoinToleranceQuantizationRegression.lean} & regression wrappers for finite-coin tolerance quantization \\
\texttt{PNP/FiniteCoinPostprocessingFoundation.lean} & exact and approximate finite-coin postprocessing laws \\
\texttt{PNP/FiniteCoinPostprocessingFoundationRegression.lean} & regression wrappers for finite-coin postprocessing laws \\
\bottomrule
\end{tabular}
\end{center}

The point of the table is simple: the grassroots route is not starting from zero.
We already have a serious computability substrate, an observer-relative semantics
layer, and a new optimistic interface for small hypothesis classes.

\section{Finite-Count Conditional Independence}

The grassroots route also needs a probability-free language for conditioning and
decorrelation.  This matters because several manuscript-style repairs move
between exact source events, randomized recodings, and output observations.
Before any asymptotic claim is meaningful, the finite counting statement must be
unambiguous.

The base exact definition is \texttt{CountIndependentWithin} from
\texttt{ConditioningObstruction.lean}.  It says that events \(E\) and \(F\) are
conditionally independent inside a finite conditioning event \(C\) by the
division-free equation
\[
  |C|\,|C\cap E\cap F| = |C\cap E|\,|C\cap F|.
\]
The approximate layer in \texttt{ApproximateDecorrelationObstruction.lean}
orients the two possible failures of this equation as a forward and backward
natural-number gap.  The new grassroots module
\texttt{FiniteCountIndependenceFoundation.lean} packages those two gaps into a
single symmetric defect:
\[
  \texttt{countIndependentWithinDefect}
  = \max(\texttt{forwardGap},\texttt{backwardGap}).
\]

\begin{theorem}[Tolerance is max-defect boundedness]
For finite events \(C,E,F\),
\[
  \texttt{CountIndependentWithinTolerance}\ C\ E\ F\ \tau
  \quad\Longleftrightarrow\quad
  \texttt{countIndependentWithinDefect}\ C\ E\ F \le \tau.
\]
\end{theorem}

\noindent
This is \texttt{countIndependentWithinTolerance\_iff\_defect\_le}.  The same
module proves monotonicity in the tolerance and the exact bridge
\[
  \texttt{CountIndependentWithinTolerance}\ C\ E\ F\ 0
  \quad\Longleftrightarrow\quad
  \texttt{CountIndependentWithin}\ C\ E\ F,
\]
formalized as
\texttt{countIndependentWithinTolerance\_zero\_iff\_countIndependentWithin}.
Equivalently,
\[
  \texttt{countIndependentWithinDefect}\ C\ E\ F = 0
\]
is exactly the original cross-multiplied finite-count independence equation.

The module also lifts this to output variables.  A certificate
\texttt{CountIndependentWithinToleranceOutput} specializes to every output
fiber, with direct projection lemmas for the forward and backward defects.  At
zero tolerance it is exactly exact finite-count independence for every output
fiber:
\[
  \texttt{countIndependentWithinToleranceOutput\_zero\_iff\_countIndependentWithin\_fibers}.
\]
For a finite output type, the file defines
\texttt{countIndependentWithinOutputDefect}, the maximum fiber defect over all
outputs, and proves
\[
  \texttt{CountIndependentWithinToleranceOutput}\ C\ E\ Y\ \tau
  \quad\Longleftrightarrow\quad
  \texttt{countIndependentWithinOutputDefect}\ C\ E\ Y \le \tau.
\]

This is intentionally small infrastructure, but it changes the shape of later
work.  Future positive routes can target a clean finite-count certificate; crux
results can be read as applications showing that certain recoding interfaces
force that certificate to have a large defect.

\section{Finite-Coin Source Erasure}

The finite-count layer above deliberately treats decorrelation as a property of
events and output fibers.  A separate grassroots question is what an output
recoding can still preserve once every finite-coin output fiber is balanced
between the true and false source sides.  The answer is now formalized in
\texttt{FiniteCoinSourceErasureFoundation.lean}.

For a recoding \(r : \mathrm{Bool} \to \mathrm{Coin} \to \alpha\), the file
defines \texttt{finiteCoinOutputPredicateNeutral}: a decidable output predicate
\(P : \alpha \to \mathrm{Prop}\) is neutral when it selects the same number of
true-side and false-side coin samples.  It then defines
\texttt{FiniteCoinOutputPredicateErasing}\(r\) to mean that every decidable
output predicate is neutral.

\begin{theorem}[Balanced fibers are predicate erasure]
If \(\alpha\) has decidable equality, then
\[
  \texttt{FiniteCoinRecodingFiberBalanced}\ r
  \quad\Longleftrightarrow\quad
  \texttt{FiniteCoinOutputPredicateErasing}\ r.
\]
\end{theorem}

\noindent
This is
\texttt{finiteCoinRecodingFiberBalanced\_iff\_outputPredicateErasing}.  The
forward direction says balanced fibers erase not only singleton outputs but every
decidable observation of the output.  The reverse direction uses singleton
predicates, so the equivalence is exact rather than an artifact of one chosen
statistic.

The same file records strict refuters.  Any output predicate with a strict
true-side or false-side count advantage contradicts predicate erasure.  If the
coin space is nonempty, a predicate that uniformly separates the two source
sides also contradicts predicate erasure.  A structural specialization is now
included as \texttt{FiniteCoinSourceRangesSeparated}: disjoint true-side and
false-side output ranges refute predicate erasure, and with decidable output
equality they also construct an explicit Boolean source decoder by testing
membership in the true-side output range.

Finally, the file packages the decoder consequence.  For a Boolean decoder
\(\mathrm{decode} : \alpha \to \mathrm{Bool}\), it defines the aggregate correct
count
\[
  \#\{c : \mathrm{Coin} \mid \mathrm{decode}(r(\mathrm{true},c))=\mathrm{true}\}
  +
  \#\{c : \mathrm{Coin} \mid \mathrm{decode}(r(\mathrm{false},c))=\mathrm{false}\}.
\]
Predicate erasure forces this count to be exactly \(|\mathrm{Coin}|\), i.e. half
of the two-sided sample space in division-free form.  Therefore any decoder with
strictly better-than-half aggregate recovery, or any decoder that uniformly
recovers the source bit on a nonempty coin space, refutes erasure.

Equivalently, the file defines \texttt{FiniteCoinOutputSourceRecovering} for the
existence of such a decoder and proves that source-recovering outputs are
incompatible with predicate erasure.  This keeps the future positive route
honest: preserving the source bit at the output level is a named mathematical
property, not a prose-side intuition.

This is a grassroots strengthening of the balanced-fiber crux.  It does not
merely say that the manuscript's current output statistic fails.  It proves that
the whole finite-coin balanced-fiber regime is source-erasing at the level of
all decidable output predicates and Boolean decoders.

\subsection{Source Erasure as Independence}

The source-erasure layer is now tied back to the finite-count independence API in
\texttt{FiniteCoinIndependenceBridge.lean}.  The file works on the concrete
two-sided sample space \(\mathrm{Bool}\times\mathrm{Coin}\).  It defines
\texttt{finiteCoinSourceTrue} for the source-bit event and
\texttt{finiteCoinOutput} for the output variable induced by a finite-coin
recoding.

The first part of the bridge is counting infrastructure:
\[
  |\{(b,c) : b=\mathrm{true}\}| = |\mathrm{Coin}|,
\]
and for every decidable output predicate \(P\),
\[
  |\{(b,c) : P(r(b,c))\}|
  =
  |\{c : P(r(\mathrm{true},c))\}|
  +
  |\{c : P(r(\mathrm{false},c))\}|.
\]

These identities support the exact theorem
\[
  \texttt{finiteCoinOutputPredicateNeutral}\ r\ P
  \quad\Longleftrightarrow\quad
  \texttt{CountIndependentWithin}\ \top\
    \texttt{finiteCoinSourceTrue}\ (P\circ \texttt{finiteCoinOutput}\ r).
\]
This is formalized as
\texttt{finiteCoinOutputPredicateNeutral\_iff\_countIndependentWithin\_trueSource\_outputPredicate}.
Consequently, predicate erasure is exactly independence from the source bit for
every decidable output predicate.

For output variables with decidable equality, the bridge specializes this to the
fiberwise API:
\[
  \texttt{FiniteCoinRecodingFiberBalanced}\ r
  \quad\Longleftrightarrow\quad
  \texttt{CountIndependentWithinToleranceOutput}\ \top\
    \texttt{finiteCoinSourceTrue}\ (\texttt{finiteCoinOutput}\ r)\ 0.
\]
For finite output types, it also proves the max-defect form:
\[
  \texttt{FiniteCoinRecodingFiberBalanced}\ r
  \quad\Longleftrightarrow\quad
  \texttt{countIndependentWithinOutputDefect}\ \top\
    \texttt{finiteCoinSourceTrue}\ (\texttt{finiteCoinOutput}\ r) = 0.
\]

This matters for the ground-up route because the balanced-fiber branch is no
longer merely a named escape hatch.  It is exactly the zero-defect source/output
independence condition in the same finite-count language used by the crux
obstructions.

\subsection{Defect Equals Scaled Fiber Imbalance}

The next quantitative step is in \texttt{FiniteCoinDefectFormula.lean}.  For an
output fiber \(y\), define its finite-coin imbalance as
\[
  \max(T_y-F_y,\;F_y-T_y),
\]
where \(T_y\) and \(F_y\) are the true-side and false-side coin-fiber counts.
The file proves that zero imbalance is exactly \(T_y=F_y\).

More importantly, it computes the oriented finite-count defects on
\(\mathrm{Bool}\times\mathrm{Coin}\):
\[
  \texttt{forwardGap}(y)=|\mathrm{Coin}|(T_y-F_y),
  \qquad
  \texttt{backwardGap}(y)=|\mathrm{Coin}|(F_y-T_y).
\]
Therefore the symmetric fiber defect is exactly
\[
  |\mathrm{Coin}| \cdot \max(T_y-F_y,\;F_y-T_y).
\]

The approximate output certificate has an exact reformulation:
\[
  \texttt{CountIndependentWithinToleranceOutput}\ \top\
    \texttt{finiteCoinSourceTrue}\ (\texttt{finiteCoinOutput}\ r)\ \tau
  \quad\Longleftrightarrow\quad
  \forall y,\;|\mathrm{Coin}|\cdot\texttt{fiberImbalance}(r,y)\le \tau.
\]
For finite output types, \texttt{countIndependentWithinOutputDefect} is the
finite supremum of these scaled fiber imbalances.

This turns the balanced-fiber story into a quantitative foundation: exact
balance is zero scaled imbalance, and approximate decorrelation is no stronger
or weaker than bounding every scaled fiber imbalance.

\subsection{Tolerance Quantization}

\texttt{FiniteCoinToleranceQuantization.lean} records the immediate but useful
threshold consequence of the scaled-defect formula.  If a finite-coin fiber
imbalance is nonzero, then its scaled contribution is at least one whole
\(|\mathrm{Coin}|\) block:
\[
  \texttt{fiberImbalance}(r,y)\ne 0
  \quad\Longrightarrow\quad
  |\mathrm{Coin}|\le
  |\mathrm{Coin}|\cdot\texttt{fiberImbalance}(r,y).
\]
This is
\texttt{finiteCoin\_card\_le\_scaled\_fiberImbalance\_of\_ne\_zero}.

Consequently, any scaled fiber-imbalance bound below \(|\mathrm{Coin}|\) forces
that fiber to be exactly balanced.  At the output-variable level, the file proves
\[
  \tau < |\mathrm{Coin}|
  \quad\Longrightarrow\quad
  \left(
  \texttt{CountIndependentWithinToleranceOutput}\ \top\
    \texttt{finiteCoinSourceTrue}\ (\texttt{finiteCoinOutput}\ r)\ \tau
  \Longleftrightarrow
  \texttt{FiniteCoinRecodingFiberBalanced}\ r
  \right).
\]
This is
\texttt{countIndependentWithinToleranceOutput\_finiteCoinOutput\_lt\_card\_iff\_fiberBalanced}.

For finite output types, the max-defect form is also pinned down.  If the
maximum output defect is below \(|\mathrm{Coin}|\), the recoding has balanced
fibers; conversely, every unbalanced finite-output recoding has output defect at
least \(|\mathrm{Coin}|\).  When the coin space is nonempty, defect below one
coin block is equivalent to balanced fibers:
\[
  \texttt{countIndependentWithinOutputDefect}\ \top\
    \texttt{finiteCoinSourceTrue}\ (\texttt{finiteCoinOutput}\ r)
  < |\mathrm{Coin}|
  \quad\Longleftrightarrow\quad
  \texttt{FiniteCoinRecodingFiberBalanced}\ r.
\]

This gives the finite-coin grassroots layer a discrete threshold theorem:
``tiny'' approximate source/output independence, below one coin block of
finite-count defect, is not a genuine relaxation of balanced-fiber erasure.

\subsection{Deterministic Postprocessing}

\texttt{FiniteCoinPostprocessingFoundation.lean} adds the postprocessing laws
needed to keep the finite-coin layer stable under ordinary recodings of the
output.  If \(g : \alpha \to \beta\) is a deterministic map, exact predicate
erasure is preserved:
\[
  \texttt{FiniteCoinOutputPredicateErasing}\ r
  \quad\Longrightarrow\quad
  \texttt{FiniteCoinOutputPredicateErasing}\ (g\circ r).
\]
With decidable equality on the source and postprocessed output types, this gives
the balanced-fiber form:
\[
  \texttt{FiniteCoinRecodingFiberBalanced}\ r
  \quad\Longrightarrow\quad
  \texttt{FiniteCoinRecodingFiberBalanced}\ (g\circ r).
\]

The approximate statement is deliberately quantitative.  For finite original
output type \(\alpha\), the file defines
\texttt{finiteOutputPreimageFiber}\((g,z)\) and
\texttt{finiteOutputMapFiberCardBound}\((g,m)\), meaning every postprocessed
fiber merges at most \(m\) original output values.  It proves the count
identities
\[
  T'_z=\sum_{y:g(y)=z} T_y,\qquad
  F'_z=\sum_{y:g(y)=z} F_y,
\]
and the imbalance inequality
\[
  \texttt{fiberImbalance}(g\circ r,z)
  \le
  \sum_{y:g(y)=z}\texttt{fiberImbalance}(r,y).
\]

Consequently, if \(r\) has finite-coin output tolerance \(\tau\), then
\[
  \texttt{CountIndependentWithinToleranceOutput}(r,\tau)
  \quad\Longrightarrow\quad
  \texttt{CountIndependentWithinToleranceOutput}(g\circ r,m\tau).
\]
This is
\texttt{countIndependentWithinToleranceOutput\_finiteCoinOutput\_postcompose}.
For finite postprocessed output types, the max-defect form is:
\[
  \texttt{outputDefect}(g\circ r)
  \le
  m\,\texttt{outputDefect}(r).
\]
The file also records the trivial bound \(m=|\alpha|\), so every finite
postprocessing has an explicit worst-case blowup.

This blocks another common ambiguity in repair attempts.  Deterministic
postprocessing cannot restore source information after exact erasure; for
approximate certificates it can only aggregate already present fiber defects,
with the aggregation controlled by finite preimage size.

\subsection{Route-Level Scoping}

The finite-coin stack is useful, but it is not a stop-grade route closure by
itself.  The crux synthesis ledger now makes this explicit in
\texttt{CruxSynthesis.lean}.  The covered finite-coin subrepair classes are:
same-source finite-coin marginal recoding, predicate erasure, decoder
half-accuracy, exact scaled fiber imbalance, sub-cardinality tolerance
quantization, deterministic exact postprocessing, and finite-preimage
approximate postprocessing with \(\tau\mapsto m\tau\) blowup.
The same synthesis ledger now also tracks a narrow finite-coin
randomized-residual subledger: supportwise unresolved product-space residuals
have no positive advantage, and any positive finite-coin randomized-residual
advantage must produce a positive-weight resolving source/coin pair.  On the
exact post-switch surface this becomes a positive-weight input/coin pair that
distinguishes the local-input switch.

The same ledger explicitly records what this does \emph{not} cover.  The broad
\texttt{randomizedResidual}, \texttt{residualSideInformation},
\texttt{approximateDecorrelation}, and \texttt{kpolyCompressionBridge} repair
classes remain outside the current packet.  Thus the finite-coin work should be
read as a precise foundation for a same-source finite-count subroute and as
route-level guidance for future attacks, not as a proof that all randomized or
approximate repairs are blocked.

Because the finite-coin subledger is now explicit, the synthesis ledger selects
\texttt{kpolyCompressionBridge} as the next highest-marginal target.  This is
the remaining manuscript-facing bridge after the abstract clocked
finite-uniform counting kernel has been separated and covered.
\texttt{ClockedKpolyBridgeEquivalence.lean} now tightens that target further:
existence of a clocked \(s\)-bit realization family is equivalent to the
existing exact visible compression target, and changing the external clock
bound does not change realizability.  So the next target is a real
small-class theorem for the switched family, not a clock-bookkeeping repair.
The constructive direction is now explicit rather than hidden inside the
existential proof.  \texttt{ClockedKpolyBridge.lean} defines
\texttt{ClockedBitCodeFamily.ofBitEncodedClassifierFamily}, which equips an
ordinary \(s\)-bit classifier family with an external clock without changing
its decoder, realization predicate, exact-recovery rate, or non-deceptive
sample rate.  The target
interface uses this constructor in
\texttt{IndexedPredictorFamily.clockedBitCodeFamilyOfHasBitBudget}; Lean now
also proves that the underlying bit family of this constructed clocked family
realizes the original indexed family and re-witnesses the same bit budget.
Finally,
\texttt{IndexedPredictorFamily.realizedByClockedBitFamily\_clockedBitCodeFamilyOfHasBitBudget}
and
\texttt{realizedByClockedExactVisibleFamily\_of\_exactVisibleCompressionTarget}
state the construction as an actual family-level clocked realization theorem.
Thus the bridge no longer depends on a nameless clocked witness: given a
bit-budget theorem, the clocked family is canonically constructed.  The
family-level predicate
\texttt{IndexedPredictorFamily.RealizedByClockedBitFamily} and its forgetful
theorem
\texttt{IndexedPredictorFamily.hasBitBudget\_of\_realizedByClockedBitFamily}
now live in \texttt{ClockedKpolyBridge.lean} itself.  Consequently
\texttt{ClockedKpolyBridgeEquivalence.lean} no longer imports the obstruction
module; the obstruction file consumes the core realization interface rather
than defining it.

The bridge now carries the full finite-success rate pair rather than only the
realized-target half.  \texttt{ClockedKpolyBridge.lean} defines
\texttt{ClockedBitCodeFamily.nondeceptiveRate} and proves the same
\(1-2^s\rho^m\) lower bound and strict positivity statements that the
underlying bit-family layer already had.  The target interface then lifts these
to any exact \(s\)-bit budget and to any exact visible compression target via
\texttt{IndexedPredictorFamily.nondeceptiveRate\_ge\_one\_sub\_uniformMissRatio\_pow\_of\_hasBitBudget}
and
\texttt{nondeceptiveRate\_ge\_one\_sub\_uniformMissRatio\_pow\_of\_exactVisibleCompressionTarget}.
This matters because the manuscript-facing bridge now exposes both
target-independent non-deceptive samples and exact recovery for realized
targets from the same finite-count hypothesis.

The bridge also lifts the pure finite-count existence layer.  Below the
bit-budget threshold,
\texttt{ClockedBitCodeFamily.exists\_nondeceptiveSample\_of\_bitBudget\_bound\_lt}
produces an explicit non-deceptive point sample for the underlying clocked
family, and
\texttt{ClockedBitCodeFamily.exists\_sample\_empiricalRiskPredictor\_eq\_target\_of\_bitBudget\_bound\_lt}
produces a point sample on which ERM exactly recovers any realized target.  The
target interface lifts both statements to exact \texttt{HasBitBudget} witnesses
and to exact visible compression targets.  Thus the clocked bridge now exposes
both quantitative rates and concrete sample witnesses, while still leaving the
small-image theorem for the manuscript's switched family as the real burden.

\texttt{ClockedKpolyGapAssessment.lean} makes this assessment explicit.  It
defines \texttt{ClockedKpolyFiniteLearningPayload}, bundling a clocked exact
visible realization family with the threshold good-sample and exact ERM
recovery witnesses for every selected predictor.  Lean proves
\[
  \texttt{ClockedKpolyFiniteLearningPayload}\ G\ s\ \texttt{clock}
  \quad\Longleftrightarrow\quad
  \texttt{ExactVisibleCompressionTarget}\ G\ s
  \quad\Longleftrightarrow\quad
  \texttt{FinitePredictorCover}\ G\ (2^s).
\]
It also lifts the existing surjectivity and injective finite-probe lower bounds
to this full payload.  So the accumulated clocked bridge is locally complete:
all finite-learning consequences now follow from the exact visible compression
target.  But the route is not closed, because the target itself is still the
unproved small finite-image theorem for the actual switched family.

\texttt{ClockedKpolyActualGapClosure.lean} resolves the next fork for the
formalized exact ZAB ERM route.  For the concrete ERM-selected family
\texttt{exactZABDecisionListERMFamily}, Lean proves the clocked finite-learning
payload at budget \(r+2k+1\):
\texttt{exactZABDecisionListERMClockedKpolyFiniteLearningPayload}.  The
bit-valued manuscript-facing specialization is
\texttt{bitVecZABDecisionListERMClockedKpolyFiniteLearningPayload}.  For an
arbitrary family variable \(G\), the minimal transfer assumption is the exact
identification
\[
  G =
  \texttt{bitVecZABDecisionListERMFamily}\ \texttt{samples},
\]
captured by
\texttt{clockedKpolyFiniteLearningPayload\_of\_eq\_bitVecZABDecisionListERMFamily};
the packaged recovery-data interface also closes the same payload via
\texttt{BitVecZABERMRecoveryData.clockedKpolyFiniteLearningPayload}.  Without
such an exact-family identification, the present interface remains
insufficient: the full exact-visible Boolean family is still a legal
\texttt{ExactVisibleSwitchedFamily}, and
\texttt{currentExactVisibleInterface\_not\_forall\_clockedKpolyFiniteLearningPayload\_of\_lt\_surfaceCard}
refutes any theorem deriving the clocked payload for every exact-visible
family below the full visible truth-table budget.  The synthesis ledger now
absorbs this exact fork in
\texttt{currentPNPStrongestHonestStatus}: the exact ZAB closure is recorded as
narrow \(Kpoly\) subrepair progress, while the broad manuscript-specific
\texttt{kpolyCompressionBridge} remains explicitly uncovered.

\texttt{ActualSwitchedLocalFamily.lean} moves the same fork one step closer to
the manuscript's actual local normal form.  It defines
\texttt{ActualSwitchedLocalInterface} with fields
\[
  z : \Phi \to Z,\qquad
  a : I \to \Phi \to \{0,1\}^k,\qquad
  b : \Phi \to \{0,1\}^k,
\]
a local rule \(h_i : (z,a,b)\to\mathrm{Bool}\), and the equation
\[
  output_i(\phi)=h_i(z(\phi),a_i(\phi),b(\phi)).
\]
For this actual-local interface, Lean proves the same conditional closure:
if the induced exact-visible predictor family is exactly
\texttt{bitVecZABDecisionListERMFamily}\(\;\texttt{samples}\), then
\texttt{ClockedKpolyFiniteLearningPayload} holds at budget \(r+2k+1\).
But Lean also proves the locality-only obstruction.  Taking blocks themselves
to be exact visible inputs and indices to be all Boolean rules gives a legal
actual-local interface whose predictor family is the full exact-visible
Boolean family.  Hence, below the full visible truth-table budget,
\texttt{actualSwitchedLocalInterface\_not\_forall\_clockedKpolyFiniteLearningPayload\_of\_lt\_surfaceCard}
refutes any derivation of the clocked payload from the current actual-local
normal-form assumptions alone.

\texttt{ActualSwitchedLocalStructure.lean} records the next narrowed proof
obligation.  The extra assumption is not more locality; it is that the local
rules selected by the actual switched family come from a bounded code class.
The structure \texttt{ActualSwitchedLocalInterface.BitCodeData} supplies one
\(s\)-bit encoded family and proves
\[
  \texttt{ExactVisibleCompressionTarget}\ T.\texttt{predictorFamily}\ s,
  \qquad
  T.\texttt{predictorFamily.FinitePredictorCover}(2^s),
\]
and therefore \texttt{ClockedKpolyFiniteLearningPayload} at the same budget.
This is the weakest formal closure currently isolated: with a finite local
code/quotient map, the small-image gap closes.  The same file proves that this
condition is substantive by refuting it for
\texttt{fullRuleActualSwitchedLocalInterface} whenever
\(s\) is below the exact visible truth-table budget.

The file also gives a concrete manuscript-shaped sufficient instance:
\texttt{ActualSwitchedLocalInterface.ZABDecisionListData}.  It requires the
actual local rules to be realized by the existing shared
\((zfeat(z),a,b)\) decision-list family, without requiring equality to a
particular ERM selector.  Lean then proves a finite predictor cover, finite
representative-index cover, finite predictor quotient, and the clocked payload
at budget \(r+2k+1\).  Under
\[
  r+2k+1 < |\texttt{ExactVisiblePostSwitchSurface}\ Z\ k|,
\]
Lean also proves that no such ZAB-structured family can have a surjective
full Boolean predictor image, so the prior full-rule local counterexample is
excluded.  Thus the current exact status is a minimal missing-assumption
result: locality-only is insufficient; bounded local code data closes the
finite-image target; and shared ZAB decision-list realization is one concrete
candidate for the manuscript-derived non-arbitrary structure that would supply
that code data.
The exact ERM equality hypothesis now factors through this narrowed interface:
\texttt{zabDecisionListData\_of\_eq\_exactZABDecisionListERMFamily} constructs
\texttt{ZABDecisionListData} from equality to the concrete shared ERM family,
and
\texttt{bitCodeData\_of\_eq\_exactZABDecisionListERMFamily} constructs the
bounded-code assumption itself.  The bit-valued identity-extractor version has
matching constructors for
\texttt{bitVecZABDecisionListERMFamily}.  Thus exact ERM equality is now a
sufficient way to discharge the minimal finite-image assumption, not a separate
parallel shortcut.

\texttt{ActualSwitchedLocalERMConstruction.lean} removes one layer of that
conditionality for the ERM-shaped construction itself.  Given concrete
manuscript-visible maps
\[
  zOf : \Phi\to Z,\qquad
  aOf : I\to\Phi\to\{0,1\}^k,\qquad
  bOf : \Phi\to\{0,1\}^k,
\]
an extractor \(zfeat : Z\to\{0,1\}^r\), and indexed samples, Lean constructs
\[
  \texttt{exactZABDecisionListERMConstruction}\ zOf\ aOf\ bOf\ zfeat\
  samples : \texttt{ActualSwitchedLocalInterface}.
\]
Its predictor family is definitionally
\texttt{exactZABDecisionListERMFamily}\((zfeat,samples)\).  Consequently Lean
proves, with no additional assumption, that this constructed actual-local
family satisfies
\texttt{ActualSwitchedLocalInterface.ZABDecisionListData} and
\texttt{ActualSwitchedLocalInterface.BitCodeData} at budget \(r+2k+1\), and
then proves
\[
  \texttt{FinitePredictorCover}\ (2^{r+2k+1})
  \quad\text{and}\quad
  \texttt{ClockedKpolyFiniteLearningPayload}\ (r+2k+1)\ clock .
\]
The bit-valued identity-extractor specialization is formalized as
\texttt{bitVecZABDecisionListERMConstruction}.  This is now a concrete
construction theorem, not merely a strengthened-interface hypothesis.  The
remaining manuscript-specific obligation is exact and narrower: prove that the
paper's switched family is this ERM construction, or else supply the same
per-index ZAB decision-list codes through
\texttt{zabDecisionListDataOfCodes}.  Without that identification, the earlier
locality-only full-rule counterexample still shows that actual-local shape
alone cannot imply a small predictor image.

\texttt{ActualSwitchedLocalCodeConstruction.lean} formalizes that second route
directly.  Given the same visible maps \(zOf,aOf,bOf\), extractor \(zfeat\),
and a function
\[
  codes : I\to
    \texttt{SharedAffineDecisionListCode}\,(r+(k+k)),
\]
Lean constructs
\texttt{exactZABCodeConstruction}.  Its exact-visible predictor family is
definitionally the canonical ZAB code family induced by those codes, and the
actual output is the corresponding decision-list prediction evaluated at
\((zOf(\phi),aOf_i(\phi),bOf(\phi))\).  From this concrete construction Lean
proves
\texttt{exactZABCodeConstruction\_zabDecisionListData},
\texttt{exactZABCodeConstruction\_bitCodeData},
\texttt{exactZABCodeConstruction\_finitePredictorCover}, and
\texttt{exactZABCodeConstruction\_clockedKpolyFiniteLearningPayload}.  Thus an
explicit per-index ZAB code assignment is now enough to close the exact
small-image target for an actual switched-local wrapper, independently of the
ERM selector.  The manuscript burden is correspondingly reduced to either an
ERM equality proof or a code extraction proof for its local rules.

\texttt{ActualSwitchedLocalPluginLookupObstruction.lean} checks the literal
finite-alphabet plug-in endpoint described in the v11/arXiv manuscript
paragraph: the local rule is implemented by a hash-table lookup on the exact
post-switch alphabet \(u=(z,a_i,b)\), with no hypothesis enumeration.  Lean
formalizes that unrestricted endpoint as
\texttt{pluginLookupActualSwitchedLocalInterface}.  Its predictor family is the
full exact-visible Boolean rule family, and
\texttt{pluginLookupActualSwitchedLocalInterface\_surjective} proves that every
Boolean lookup table on the post-switch alphabet is realized.  Consequently,
below the exact visible truth-table budget, Lean proves
\[
  \texttt{pluginLookupActualSwitchedLocalInterface\_not\_finitePredictorCover\_of\_lt\_surfaceCard},
\]
the corresponding clocked-payload obstruction, and the two constructor
obstructions
\texttt{pluginLookupActualSwitchedLocalInterface\_not\_bitCodeData\_of\_lt\_surfaceCard}
and
\texttt{pluginLookupActualSwitchedLocalInterface\_not\_zabDecisionListData\_of\_lt\_surfaceCard}.
Thus a finite-alphabet lookup implementation is not itself the missing
small-image theorem.  It closes only after an additional manuscript theorem
identifies those lookup tables with the ZAB ERM family or with explicit
per-index ZAB codes.

\texttt{ActualSwitchedLocalPluginMajorityObstruction.lean} removes a possible
escape hatch in that diagnosis.  The manuscript states the local rule as a
plug-in majority over training labels grouped by the visible input \(u\).  Lean
formalizes the aggregated version as \texttt{PluginMajorityCounts}: for each
visible input, record the number of true and false votes and return the
majority label.  The theorem
\texttt{pluginMajorityActualSwitchedLocalInterface\_realizes\_every\_lookupTable}
proves that every Boolean lookup table is still realized, by placing one vote
on the desired side for each visible input.  Therefore
\texttt{pluginMajorityActualSwitchedLocalInterface\_surjective} and the
matching no-cover/no-payload/no-\texttt{BitCodeData}/no-\texttt{ZABDecisionListData}
theorems show that the counted majority endpoint is also full-rule expressive
below the exact visible truth-table budget.  The obstruction is therefore not
caused by replacing majority with a lookup table; it is caused by the absence
of a proved restriction from the plug-in majority counts to the small ZAB/code
class.

\texttt{ActualSwitchedLocalPluginSampleMajorityObstruction.lean} removes the
last aggregation-only caveat from the plug-in majority blocker.  Instead of
postulating the true/false vote counts directly, Lean defines
\texttt{PluginSampleMajority.trueVotes} and
\texttt{PluginSampleMajority.falseVotes} by counting labels in an actual finite
sample.  For every exact-visible Boolean rule \(rule\), the graph sample
\[
  \{(u,rule(u)) : u\in \texttt{ExactVisiblePostSwitchSurface}\,Z\,k\}
\]
has one labeled example at each visible input.  Lean proves
\texttt{PluginSampleMajority.predict\_graphSample}, and then
\texttt{pluginSampleMajorityActualSwitchedLocalInterface\_surjective}: the
sample-level plug-in majority interface is still surjective onto all Boolean
rules on \(u=(z,a,b)\).  The companion no-cover/no-payload/no-\texttt{BitCodeData}/
no-\texttt{ZABDecisionListData} theorems therefore show that finite samples and
empirical majority alone do not imply the manuscript's needed small predictor
image.  A successful route must still prove that the actual samples/rules are
restricted to the existing exact ZAB ERM family or to explicit bounded ZAB
codes.

\texttt{ActualSwitchedLocalBoundedSampleMajorityObstruction.lean} checks the
next natural repair: bounding the finite sample length.  Lean proves
\texttt{PluginSampleMajority.length\_graphSample}, namely that the graph sample
of a rule has length exactly the cardinality of the visible alphabet.  It then
defines a bounded-sample plug-in majority interface whose indices are samples
of length at most \(B\).  If
\[
  |\texttt{ExactVisiblePostSwitchSurface}\,Z\,k|\le B,
\]
Lean proves
\texttt{boundedSamplePluginMajorityActualSwitchedLocalInterface\_surjective}.
The strict opposite case is now also machine checked.  Lean proves
\texttt{boundedSamplePluginMajorityActualSwitchedLocalInterface\_not\_surjective\_of\_sampleBound\_lt\_surfaceCard}:
if \(B<|\texttt{ExactVisiblePostSwitchSurface}\,Z\,k|\), every bounded sample
omits some visible input, where the current tie-breaking plug-in majority
predicts \texttt{true}; hence the all-false rule is not realizable.  Combining
the two directions gives
\texttt{boundedSamplePluginMajorityActualSwitchedLocalInterface\_surjective\_iff}.
Consequently, sample length alone has an exact full-rule expressivity threshold
at the visible alphabet size.  Above that threshold the endpoint still has no
finite predictor cover, clocked payload, \texttt{BitCodeData}, or
\texttt{ZABDecisionListData} below the exact visible truth-table budget.  Below
that threshold Lean proves only non-surjectivity; a separate image-size,
quotient, ZAB ERM, or code theorem is still needed to close the manuscript's
small-image obligation.

\texttt{ExactVisibleImageBudget.lean} refines the same target in grassroots
terms.  Lean proves that an exact \(s\)-bit budget is equivalent to a finite
cover of the predictor image by at most \(2^s\) Boolean functions, and proves
the same equivalence for the clocked exact-visible interface.  So the remaining
\texttt{kpolyCompressionBridge} burden can now be phrased without implementation
noise: prove that the actual switched family has a finite predictor image of
the required size.  The theorem does not assert that this image is small; it
identifies the exact finite-image obligation any successful repair must meet.
The same module also proves that a finite predictor cover is equivalent to a
finite representative-index cover: finitely many actual selected indices whose
predictors represent every predictor in the family.  This avoids treating the
cover as an arbitrary external set of Boolean functions.
It also proves the equivalent finite quotient-code presentation: every index is
encoded into one of finitely many code values, and the decoded Boolean predictor
is exactly the selected predictor for that index.  The clocked exact-visible
interface now has matching negative forms too: nonexistence of a clocked
realization is equivalent to nonexistence of each of these finite-image
presentations at the same \(2^s\) budget.
It now includes the transport layer needed for later route-building:
predictor-image covers, representative-index covers, and quotient-code
presentations are monotone in the allowed budget, pull back along an explicit
factorization through a visible summary, and push back to that summary when the
view has a section.  This makes the finite-image obligation stable under the
same kind of exact pullback arguments used elsewhere in the PNP development,
without claiming that the actual switched family has a small image.
The reduced-view lower bound is now also machine checked.  If an exact family
factors through a finite visible summary with a section, and the reduced family
is surjective onto all Boolean classifiers on that summary, then every finite
predictor-image cover, finite representative-index cover, and finite
quotient-code presentation of the exact family must have size at least
\(2^{|\mathrm{summary}|}\).  The raw \((a,b)\)-surface specialization gives
the concrete threshold \(2^{2^{2k}}\) for \(\texttt{ABVisibleSurface}\ k\).
Thus finite-image compression after quotienting to \((a,b)\) still requires an
actual strict-subclass theorem for the reduced switched predictors.
The same file now proves the matching finite-image counting obstruction: if an
indexed family on a finite surface is still surjective onto all Boolean
classifiers, then every finite predictor cover, finite representative-index
cover, and finite quotient-code presentation has size at least the full
classifier-space cardinality \(2^{|\mathrm{surface}|}\).  In the exact visible
clocked interface, this gives a direct no-go theorem below that predictor-card
threshold.
\texttt{ProductSuccessKpolyBridgeObstruction.lean} adds the corresponding
route-level sanity check.  A finite tower product bound, or even a valid
fixed-field v13 switching certificate, is not a finite-image theorem for the
exact visible switched family.  Lean refutes product-only and fielded-only
compression bridges using the full exact-visible Boolean family, including a
nonempty one-step Boolean-square certificate with the one-cell field.  So the
grassroots \(Kpoly\) target remains a semantic predictor-image theorem, not a
consequence of product success alone.
The same module now pins the positive boundary too.  Once the product premise
or the fielded switching certificate is actually true, an implication from
that premise to exact-visible compression is equivalent to a finite
predictor-image cover of size at most \(2^s\).  The global product-only and
fielded-only bridge schemas are therefore exactly schemas demanding such covers
for every exact-visible family; when instantiated at one family, they also
yield finite representative-index covers and quotient-code presentations.
This keeps the next grassroots obligation concrete: build the finite image of
the actual switched predictors, or exhibit a fully expressive subcase that
blocks it.
\texttt{CruxSynthesis.lean} now records this status in the narrow \(Kpoly\)
subrepair ledger: product-bound and fielded-switching bridge schemas are covered
only as finite-image-boundary reductions and full-visible refuters.  The broad
manuscript-facing \(Kpoly\) bridge remains a live gap until the actual switched
predictor class is proved genuinely small.
The route check is no longer tied to the canonical full-rule family: any
exact-visible family with a surjective prediction map now blocks the
product-only or fielded-only bridge below the same surface threshold, and blocks
the semantic finite-image bridge below the corresponding predictor-cardinality
threshold.  So changing the index presentation of a fully expressive family is
not a repair.

\section{Encoded Hypothesis Families}

The key new grassroots abstraction is that a hypothesis class should be represented
by a finite code space together with a decoder.

\begin{quote}
code budget first, realized class second.
\end{quote}

Formally, \texttt{HypothesisClass.lean} introduces an \emph{encoded family}
\[
  \mathrm{decode} : \mathrm{Code} \to (\mathrm{Input} \to \mathrm{Output}),
\]
where \(\mathrm{Code}\) is finite.  The realized hypothesis class is the range of
\(\mathrm{decode}\).

\begin{theorem}[Realized class is bounded by code space]
For any encoded family \(H\),
\[
  |\mathrm{realized}(H)| \le |\mathrm{Code}(H)|.
\]
\end{theorem}

\noindent
This is the theorem \texttt{EncodedFamily.card\_realized\_le} in Lean.

The specialization relevant to the present route is the bit-coded case.

\begin{theorem}[Bit-budget bound]
If a family of Boolean classifiers is decoded from \(s\) bits, then it realizes at
most \(2^s\) distinct classifiers.
\end{theorem}

\noindent
This is formalized as
\texttt{BitEncodedClassifierFamily.card\_realized\_le}.

This theorem is elementary, but it captures the exact shape of the optimistic
repair that would be needed later.  A switching theorem cannot stop at ``the new
predictor is local.''  It must continue to a statement of the form:

\begin{quote}
the switched predictors lie in a family decoded from \(\mathrm{polylog}(m)\) bits.
\end{quote}

Only then does finite-class learning theory become genuinely available.

\section{Finite-Class ERM Kernel}

The next honest step after a finite encoded family is not yet a PAC theorem.
It is the finite combinatorial kernel of empirical risk minimization.

The new module \texttt{FiniteERM.lean} defines labeled samples as finite lists of
pairs \((x,y)\), empirical error as the number of misclassified examples in the
sample, and exact sample consistency as zero empirical error.

\begin{theorem}[Finite ERM existence]
Let \(H\) be a nonempty encoded family.  For every finite labeled sample \(S\),
there exists a realized predictor \(h^\star \in \mathrm{realized}(H)\) such that
\[
  \mathrm{err}_S(h^\star) \le \mathrm{err}_S(g)
  \qquad
  \text{for all } g \in \mathrm{realized}(H).
\]
\end{theorem}

\noindent
This is formalized as
\texttt{EncodedFamily.exists\_realized\_empiricalRiskMinimizer}.

\begin{theorem}[Zero-error transfer through ERM]
If some code in the family fits the sample exactly, then the chosen ERM
predictor also fits the sample exactly.
\end{theorem}

\noindent
This is formalized as
\texttt{EncodedFamily.empiricalRiskPredictor\_fitsSample\_of\_exists\_code\_fits}.

This is still modest, but it matters.  Combined with the code-budget theorem
from \texttt{HypothesisClass.lean}, it gives the first clean bridge from
``small encoded family'' to an actual learning-theoretic construction.

\section{Finite Uniform Consistency Bound}

The next step now formalized is a fully combinatorial finite-domain bound.  We
do not yet assume an arbitrary distribution or invoke measure theory.  Instead,
we count how many length-\(m\) input samples can make a wrong predictor look
perfectly consistent with the target labels.

\begin{theorem}[Wrong predictors fit only exponentially few samples]
Let \(X\) be a finite input domain.  If \(h \neq f\), then the number of
length-\(m\) samples on which \(h\) agrees with the target \(f\) at every sampled
point is at most
\[
  (|X|-1)^m.
\]
\end{theorem}

\noindent
This is formalized as
\texttt{card\_consistentSamples\_le\_of\_exists\_disagreement} in
\texttt{FiniteConsistencyBound.lean}.

\begin{theorem}[Finite encoded families have few deceptive samples]
Let \(H\) be an encoded family over a finite input domain \(X\).  The number of
length-\(m\) samples on which some wrong code in \(H\) still matches all target
labels is at most
\[
  |\mathrm{Code}(H)| \cdot (|X|-1)^m.
\]
\end{theorem}

\noindent
This is formalized as
\texttt{EncodedFamily.card\_deceptiveSamples\_le}.

\begin{theorem}[ERM recovers the target off the deceptive set]
If the target predictor is realized by the encoded family, then on every
non-deceptive sample the chosen ERM predictor equals the target predictor
exactly.
\end{theorem}

\noindent
This is formalized as
\texttt{EncodedFamily.empiricalRiskPredictor\_eq\_target\_of\_not\_deceptive}.

This is not yet the full finite-class generalization theorem one would state in
terms of \(\log |H|\) and an arbitrary data distribution.  But it is already the
right union-bound skeleton in a completely concrete finite setting.

\section{Finite Uniform Recovery Threshold}

The next corollary is the cleanest possible extraction from the deceptive-sample
bound.  If the deceptive set is strictly smaller than the whole length-\(m\)
sample space, then at least one non-deceptive sample exists.

\begin{theorem}[Threshold for existence of a good sample]
Let \(H\) be an encoded family over a finite input domain \(X\).  If
\[
  |\mathrm{Code}(H)| \cdot (|X|-1)^m < |X|^m,
\]
then there exists a length-\(m\) sample outside the deceptive set.
\end{theorem}

\noindent
This is formalized as
\texttt{EncodedFamily.exists\_nondeceptiveSample\_of\_bound\_lt} in
\texttt{FiniteUniformRecovery.lean}.

\begin{theorem}[Threshold corollary for exact ERM recovery]
If the target predictor is realized by the encoded family and
\[
  |\mathrm{Code}(H)| \cdot (|X|-1)^m < |X|^m,
\]
then there exists a length-\(m\) sample on which the chosen ERM predictor equals
the target predictor exactly.
\end{theorem}

\noindent
This is formalized as
\texttt{EncodedFamily.exists\_sample\_empiricalRiskPredictor\_eq\_target\_of\_bound\_lt}.

The same counting threshold now has an explicit bit-budget form.  For an
\(s\)-bit Boolean family \(F\), the generic code-cardinality factor becomes
\(2^s\).  Thus
\[
  2^s (|X|-1)^m < |X|^m
\]
implies the existence of a non-deceptive sample for \(F\).  If the target is
realized by \(F\), the bit-family ERM wrapper exactly recovers the target on
some labeled point sample.  These statements are formalized as
\texttt{BitEncodedClassifierFamily.exists\_nondeceptiveSample\_of\_bitBudget\_bound\_lt}
and
\texttt{BitEncodedClassifierFamily.exists\_sample\_empiricalRiskPredictor\_eq\_target\_of\_bitBudget\_bound\_lt},
with regression wrappers in
\texttt{FiniteUniformRecoveryRegression.lean}.

There is also a two-point-domain specialization.  If \(|X|=2\), the same
threshold simplifies to
\[
  s < m.
\]
Indeed \(2^s(|X|-1)^m < |X|^m\) becomes \(2^s < 2^m\).  The module
\texttt{BitFamilyUniformRecoveryBinary.lean} formalizes this as
\texttt{BitEncodedClassifierFamily.bitBudget\_bound\_lt\_of\_card\_eq\_two\_of\_lt\_sampleSize},
then derives non-deceptive-sample and exact ERM recovery corollaries for any
finite two-point input domain, with Boolean-domain wrappers as the canonical
binary case.

The clocked \(Kpoly\) bridge now carries the same existential layer.  Theorems
\texttt{ClockedBitCodeFamily.exists\_nondeceptiveSample\_of\_bitBudget\_bound\_lt}
and
\texttt{ClockedBitCodeFamily.exists\_sample\_empiricalRiskPredictor\_eq\_target\_of\_bitBudget\_bound\_lt}
lift the generic bit-budget threshold to clocked families.  On the Boolean
domain, \texttt{ClockedBitCodeFamily.exists\_bool\_nondeceptiveSample\_of\_lt\_sampleSize}
and
\texttt{ClockedBitCodeFamily.exists\_bool\_sample\_empiricalRiskPredictor\_eq\_target\_of\_lt\_sampleSize}
give the concrete \(s<m\) versions.  These are constructive consequences of a
small clocked bit family; they do not assert that the exact switched family has
one.

This is still finite and non-asymptotic, but it clarifies the next job.  The
rate layer below turns the combinatorial ratio behind the threshold into an
explicit uniform success bound; the later weighted layer relaxes the uniform
sampling assumption.

\section{Finite Uniform Rate Layer}

The next natural refinement is still finite, but no longer purely existential.
Instead of asking only whether some good sample exists, we measure how large the
good and bad regions are inside the finite sample space.

The module \texttt{FiniteUniformRate.lean} defines rational uniform densities on
the point-sample space:
\begin{enumerate}[leftmargin=1.5em]
  \item the deceptive rate,
  \item the non-deceptive rate,
  \item the exact-recovery rate for ERM.
\end{enumerate}

\begin{theorem}[Uniform deceptive-rate bound]
Let \(H\) be an encoded family over a nonempty finite input domain \(X\).  Then
the uniform fraction of length-\(m\) samples that are deceptive is at most
\[
  \frac{|\mathrm{Code}(H)| \cdot (|X|-1)^m}{|X|^m}.
\]
\end{theorem}

\noindent
This is formalized as
\texttt{EncodedFamily.deceptiveRate\_le\_bound} in
\texttt{FiniteUniformRate.lean}.

\begin{theorem}[Exact recovery dominates the non-deceptive region]
If the target predictor is realized by the encoded family, then the uniform exact
recovery rate of the chosen ERM predictor is at least
\[
  1 - \mathrm{rate}_{\mathrm{deceptive}}.
\]
\end{theorem}

\noindent
This is formalized as
\texttt{EncodedFamily.exactRecoveryRate\_ge\_one\_sub\_deceptiveRate}.

\begin{theorem}[Positive exact recovery rate under the threshold]
If the target predictor is realized and
\[
  |\mathrm{Code}(H)| \cdot (|X|-1)^m < |X|^m,
\]
then the exact recovery rate is strictly positive.
\end{theorem}

\noindent
This is formalized as
\texttt{EncodedFamily.exactRecoveryRate\_pos\_of\_bound\_lt}.

This is still not a full PAC-style theorem under arbitrary distributions.  But it
does close the uniform finite-sampling gap between ``some good sample exists'' and
``good samples occupy an explicitly bounded fraction of the sample space.''

\section{Finite Uniform Success Bounds}

The next cleanup step is to express the finite uniform rate bound using the
single-step miss ratio
\[
  \rho_X := \frac{|X|-1}{|X|}.
\]
For any nonempty finite input domain, \(\rho_X < 1\), so the finite deceptive
rate decays exponentially in the sample length \(m\).  The Lean layer now also
records the elementary order facts behind this slogan:
\[
  0 \le \rho_X \le 1,\qquad
  m \le n \Rightarrow \rho_X^n \le \rho_X^m.
\]
These are formalized as
\texttt{uniformMissRatio\_nonneg}, \texttt{uniformMissRatio\_le\_one}, and
\texttt{uniformMissRatio\_pow\_le\_pow\_of\_le}.  The resulting finite-code
failure term is antitone in the sample length:
\[
  m \le n \Rightarrow
  |\mathrm{Code}(H)|\rho_X^n \le |\mathrm{Code}(H)|\rho_X^m.
\]
This is formalized as
\texttt{EncodedFamily.codeMul\_uniformMissRatio\_pow\_le\_of\_le}, with the
bit-budget specialization
\texttt{BitEncodedClassifierFamily.bitBudget\_uniformMissRatio\_pow\_le\_of\_le}.

The module \texttt{FiniteUniformSuccess.lean} packages the previous rate theorem
into this cleaner interface.

\begin{theorem}[Miss-ratio deceptive bound]
Let \(H\) be an encoded family over a nonempty finite input domain \(X\).  Then
\[
  \mathrm{rate}_{\mathrm{deceptive}}
  \le
  |\mathrm{Code}(H)| \cdot \rho_X^m.
\]
\end{theorem}

\noindent
This is formalized as
\texttt{EncodedFamily.deceptiveRate\_le\_codeMul\_uniformMissRatio\_pow}.

\begin{theorem}[Uniform non-deceptive-rate lower bound]
For the same encoded family,
\[
  \mathrm{rate}_{\mathrm{nondeceptive}}
  \ge
  1 - |\mathrm{Code}(H)| \cdot \rho_X^m.
\]
\end{theorem}

\noindent
This is formalized as
\texttt{EncodedFamily.nondeceptiveRate\_ge\_one\_sub\_codeMul\_uniformMissRatio\_pow},
with epsilon and strict-positivity corollaries.  Unlike exact recovery, this
statement does not require the target to be realized by the family; it is the
pure finite complement of the deceptive-sample upper bound.  The monotone form
\texttt{EncodedFamily.nondeceptiveRate\_ge\_one\_sub\_codeMul\_uniformMissRatio\_pow\_of\_le}
also says that if \(m \le n\), then the weaker \(m\)-sample lower bound remains
valid for the \(n\)-sample non-deceptive rate.

\begin{theorem}[Uniform exact recovery lower bound]
If the target predictor is realized by the encoded family, then
\[
  \mathrm{rate}_{\mathrm{exact\ recovery}}
  \ge
  1 - |\mathrm{Code}(H)| \cdot \rho_X^m.
\]
\end{theorem}

\noindent
This is formalized as
\texttt{EncodedFamily.exactRecoveryRate\_ge\_one\_sub\_codeMul\_uniformMissRatio\_pow}.
Its monotone companion
\texttt{EncodedFamily.exactRecoveryRate\_ge\_one\_sub\_codeMul\_uniformMissRatio\_pow\_of\_le}
gives the same sample-extension principle for exact recovery whenever the
target is realized by the encoded family.

\begin{theorem}[\(\varepsilon\)-style corollary]
If
\[
  |\mathrm{Code}(H)| \cdot \rho_X^m \le \varepsilon,
\]
then the exact recovery rate is at least \(1 - \varepsilon\).
\end{theorem}

\noindent
This is formalized as
\texttt{EncodedFamily.exactRecoveryRate\_ge\_one\_sub\_of\_codeMul\_uniformMissRatio\_pow\_le}.

The same module now specializes these bounds to explicitly bit-encoded Boolean
families.  If \(F\) is encoded by \(s\) bits, then \(|\mathrm{Code}(F)| = 2^s\),
so the uniform deceptive rate is bounded by \(2^s\rho_X^m\), the non-deceptive
rate is at least \(1-2^s\rho_X^m\), and if the target is realized by \(F\), the
exact ERM recovery rate is at least
\[
  1 - 2^s \rho_X^m.
\]
This is formalized as
\texttt{BitEncodedClassifierFamily.deceptiveRate\_le\_bitBudget\_uniformMissRatio\_pow}
and
\texttt{BitEncodedClassifierFamily.nondeceptiveRate\_ge\_one\_sub\_bitBudget\_uniformMissRatio\_pow}
and
\texttt{BitEncodedClassifierFamily.exactRecoveryRate\_ge\_one\_sub\_bitBudget\_uniformMissRatio\_pow},
with epsilon and positivity corollaries plus regression wrappers.  The point is
not a P versus NP lower bound by itself; it is the exact finite bridge any
grassroots compression route must instantiate once it claims an explicit
\(s\)-bit predictor family.

The module also includes
\texttt{BitEncodedClassifierFamily.exactRecoveryRate\_pos\_of\_bitBudget\_bound\_lt},
which keeps the strict-positivity conclusion directly tied to the natural
cardinality threshold \(2^s(|X|-1)^m<|X|^m\).  Specializing this to the
two-point threshold produces \texttt{BitFamilyUniformSuccessBinary.lean}: on any
finite two-point input domain, and in particular on \texttt{Bool}, the simple
condition \(s<m\) implies a strictly positive uniform exact-recovery rate for
every target realized by the \(s\)-bit family.

The Boolean specialization is now quantitative as well.  Since
\[
  \rho_{\texttt{Bool}} = \frac{1}{2},
\]
the same module proves
\[
  \mathrm{rate}_{\mathrm{deceptive}}
  \le 2^s\left(\frac12\right)^m
\]
and
\[
  \mathrm{rate}_{\mathrm{nondeceptive}}
  \ge 1 - 2^s\left(\frac12\right)^m
\]
and, for every target realized by the \(s\)-bit family,
\[
  \mathrm{rate}_{\mathrm{exact\ recovery}}
  \ge 1 - 2^s\left(\frac12\right)^m.
\]
It also records the corresponding epsilon and positivity corollaries.  Thus the
binary layer now has both the simple threshold \(s<m\) and the explicit uniform
failure term \(2^s(1/2)^m\).

The same result is now also packaged in sample-margin form.  If the sample
length satisfies
\[
  s + k \le m,
\]
then
\[
  2^s\left(\frac12\right)^m \le \left(\frac12\right)^k,
\]
and therefore
\[
  \mathrm{rate}_{\mathrm{nondeceptive}}
  \ge 1 - \left(\frac12\right)^k.
\]
For every realized target, the exact ERM recovery rate also satisfies
\[
  \mathrm{rate}_{\mathrm{exact\ recovery}}
  \ge 1 - \left(\frac12\right)^k.
\]
This is formalized as
\texttt{BitEncodedClassifierFamily.bool\_bitBudget\_pow\_le\_margin\_pow} and
\texttt{BitEncodedClassifierFamily.nondeceptiveRate\_ge\_one\_sub\_bool\_margin\_pow}
and
\texttt{BitEncodedClassifierFamily.exactRecoveryRate\_ge\_one\_sub\_bool\_margin\_pow}.
The theorem is often the most convenient downstream form: every extra Boolean
sample beyond the bit budget halves the remaining uniform failure allowance.
There is also an epsilon form: if \((1/2)^k \le \varepsilon\), then under the
same margin hypothesis
\[
  \mathrm{rate}_{\mathrm{nondeceptive}} \ge 1 - \varepsilon
  \quad\text{and, for realized targets,}\quad
  \mathrm{rate}_{\mathrm{exact\ recovery}} \ge 1 - \varepsilon.
\]
This is formalized as
\texttt{BitEncodedClassifierFamily.bool\_bitBudget\_pow\_le\_of\_margin\_pow\_le}
and
\texttt{BitEncodedClassifierFamily.nondeceptiveRate\_ge\_one\_sub\_of\_bool\_margin\_pow\_le}
and
\texttt{BitEncodedClassifierFamily.exactRecoveryRate\_ge\_one\_sub\_of\_bool\_margin\_pow\_le}.
The strict positive-margin form is now proved too.  If \(0<k\) and
\(s+k\le m\), then
\[
  2^s\left(\frac12\right)^m < 1.
\]
Consequently the non-deceptive uniform rate is strictly positive, and for any
realized target the exact ERM recovery rate is strictly positive.  In Lean
these are
\texttt{BitEncodedClassifierFamily.bool\_bitBudget\_pow\_lt\_one\_of\_positive\_margin},
\texttt{BitEncodedClassifierFamily.nondeceptiveRate\_pos\_bool\_of\_positive\_margin},
and
\texttt{BitEncodedClassifierFamily.exactRecoveryRate\_pos\_bool\_of\_positive\_margin}.
The concrete \(s<m\) case is also recorded as
\texttt{BitEncodedClassifierFamily.bool\_bitBudget\_pow\_lt\_one\_of\_lt\_sampleSize}
and
\texttt{BitEncodedClassifierFamily.nondeceptiveRate\_pos\_bool\_of\_lt\_sampleSize}.

The clocked \(Kpoly\) counting bridge has the corresponding strict
positive-margin success statements.  A clocked \(s\)-bit Boolean family with
\(0<k\) and \(s+k\le m\) has strictly positive non-deceptive rate for every
target; if the target is realized by the family, it also has strictly positive
exact ERM recovery rate.  These are formalized as
\texttt{ClockedBitCodeFamily.nondeceptiveRate\_pos\_bool\_of\_positive\_margin}
and
\texttt{ClockedBitCodeFamily.exactRecoveryRate\_pos\_bool\_of\_positive\_margin}.
The concrete \(s<m\) existence theorems are
\texttt{ClockedBitCodeFamily.exists\_bool\_nondeceptiveSample\_of\_lt\_sampleSize}
and
\texttt{ClockedBitCodeFamily.exists\_bool\_sample\_empiricalRiskPredictor\_eq\_target\_of\_lt\_sampleSize}.
The recovery theorem still consumes a realized clocked bit-family; it does not
assert that the manuscript's switched decoder has such a small realization.

This is still uniform finite-sample learning rather than full distributional
generalization.  But it is the first version of the grassroots story with an
explicit success lower bound of the form ``one minus an exponentially decaying
failure term.''

\section{Finite Weighted Agreement and Family Bounds}

The next honest move is to relax the uniform-sampling assumption while still
staying fully finite.  Instead of sampling uniformly from a finite input domain,
we let one finite probability mass function \(\mu\) weight the inputs.

The module \texttt{FiniteIIDAgreement.lean} handles the one-predictor case.  For
one target \(f\) and one predictor \(h\), it defines:
\begin{enumerate}[leftmargin=1.5em]
  \item the one-step agreement mass of \(h\) with \(f\) under \(\mu\),
  \item the product mass of a length-\(m\) sample under independent draws from \(\mu\),
  \item the total mass of samples on which \(h\) agrees with \(f\) everywhere.
\end{enumerate}

\begin{theorem}[Exact weighted agreement law]
For any finite input type and finite PMF \(\mu\),
\[
  \mathrm{mass}_{\mu}(\text{all-$m$-point samples consistent with } h)
  =
  \bigl(\mathrm{agreementMass}_{\mu}(f,h)\bigr)^m.
\]
\end{theorem}

\noindent
This is formalized as
\texttt{consistentSampleMass\_eq\_agreementMass\_pow}.

\begin{theorem}[Weighted decay corollary]
If the one-step agreement mass is at most \(q\), then the total mass of
length-\(m\) consistent samples is at most \(q^m\).
\end{theorem}

\noindent
This is formalized as
\texttt{consistentSampleMass\_le\_pow\_of\_agreementMass\_le}.

The new module \texttt{FiniteIIDFamilyBound.lean} lifts this to the encoded-family
level.

\begin{theorem}[Weighted deceptive-family union bound]
Let \(H\) be a finite encoded family and \(\mu\) a finite PMF on the input type.
Then the total \(\mu^m\)-mass of deceptive samples is at most
\[
  \sum_{c \in \mathrm{BadCodes}(H,f)}
    \bigl(\mathrm{agreementMass}_{\mu}(f,\mathrm{decode}(c))\bigr)^m.
\]
\end{theorem}

\noindent
This is formalized as
\texttt{EncodedFamily.deceptiveSampleMass\_le\_badCodeAgreementMassSum}.

\begin{theorem}[Uniform weighted family corollary]
If every bad code in the family has one-step agreement mass at most \(q\), then
the total deceptive mass is at most
\[
  |\mathrm{Code}(H)| \cdot q^m.
\]
\end{theorem}

\noindent
This is formalized as
\texttt{EncodedFamily.deceptiveSampleMass\_le\_codeCard\_mul\_pow\_of\_agreementMass\_le}.

This is the first genuinely family-level finite distributional theorem in the
grassroots development.  The next real probabilistic step is no longer a union
bound itself, but a recovery or generalization theorem built on top of this
weighted family bound.

\section{Finite Weighted Recovery Bounds}

The module \texttt{FiniteIIDRecovery.lean} closes the next natural loop.  It
turns the weighted deceptive-mass bound into a weighted success statement for the
chosen ERM predictor.

\begin{theorem}[Weighted exact recovery dominates the nondeceptive mass]
If the target predictor is realized by the encoded family, then the total
\(\mu^m\)-mass of exact ERM recovery is at least
\[
  1 - \mathrm{deceptiveMass}_{\mu}.
\]
\end{theorem}

\noindent
This is formalized as
\texttt{EncodedFamily.exactRecoverySampleMass\_ge\_one\_sub\_deceptiveSampleMass}.

\begin{theorem}[Weighted exact recovery from bad-code agreement masses]
If the target predictor is realized by the encoded family, then the exact
recovery mass is at least
\[
  1 - \sum_{c \in \mathrm{BadCodes}(H,f)}
    \bigl(\mathrm{agreementMass}_{\mu}(f,\mathrm{decode}(c))\bigr)^m.
\]
\end{theorem}

\noindent
This is formalized as
\texttt{EncodedFamily.exactRecoverySampleMass\_ge\_one\_sub\_badCodeAgreementMassSum}.

\begin{theorem}[Uniform weighted exact recovery corollary]
If every bad code has one-step agreement mass at most \(q\), then the exact
recovery mass is at least
\[
  1 - |\mathrm{Code}(H)| \cdot q^m.
\]
\end{theorem}

\noindent
This is formalized as
\texttt{EncodedFamily.exactRecoverySampleMass\_ge\_one\_sub\_codeCard\_mul\_pow\_of\_agreementMass\_le}.

At this point the grassroots development has reached a fully finite,
distribution-weighted recovery theorem for ERM over encoded families.  The next
step is to lift this from exact recovery on a finite sample space to a genuine
finite-class generalization statement.

\section{What Still Counts as the Same Route}

At this stage the grassroots theory can also say something more conceptual, but
in a machine-checked way.  If a proposed rescue still wants to be \emph{the same
route} as the current switching/ERM program, then it cannot merely gesture at
weakness, GSLTs, or observer-relative semantics.  It must still instantiate the
existing finite learning interface.

The new module \texttt{SameRouteInterface.lean} packages exactly that point.  A
\emph{same-route family} is just:
\begin{enumerate}[leftmargin=1.5em]
  \item one available input surface \(\mathrm{AvailableInput}(P,B,G)\),
  \item one \(s\)-bit encoded boolean predictor family on that surface,
  \item one target predictor realized by some code in that family.
\end{enumerate}

\begin{theorem}[Same-route weighted recovery certificate]
Let \(F\) be an \(s\)-bit encoded classifier family and suppose the target
predictor is realized by some code in \(F\).  If every bad code has one-step
agreement mass at most \(q\), then the exact ERM recovery mass is at least
\[
  1 - 2^s q^m.
\]
\end{theorem}

\noindent
This is formalized as
\texttt{BitEncodedClassifierFamily.exactRecoverySampleMass\_ge\_one\_sub\_bitBudget\_mul\_pow\_of\_agreementMass\_le}.

The companion module \texttt{BitFamilyWeightedRecovery.lean} packages the same
fact in the form needed by downstream route interfaces: if the explicit failure
term \(2^s q^m\) is at most a chosen tolerance \(\varepsilon\), then exact ERM
recovery has sample mass at least \(1-\varepsilon\).  In Lean this is
\texttt{BitEncodedClassifierFamily.exactRecoverySampleMass\_ge\_one\_sub\_of\_bitBudget\_mul\_pow\_le},
with a regression wrapper in
\texttt{BitFamilyWeightedRecoveryRegression.lean}.

The point is not that this theorem proves the current switching step.  The
point is that it now marks the boundary of the current route.  A weakness/GSLT
idea that still intends to be a \emph{repair} of the present program must hand
over this sort of certificate: an explicit visible input surface, an explicit
small code space, and explicit agreement-mass control for the bad codes.  If it
does not, then it is no longer filling a gap in the current route.  It is
starting a different route.

The agreement-mass hypothesis is not a bookkeeping side condition.  The new
module \texttt{BadCodeAgreementObstruction.lean} proves the support-level test:
if a bad code agrees with the target on every positive-mass input, its
agreement mass is exactly \(1\).  Therefore no strict \(q<1\) recovery package
can exist unless each bad code disagrees with the target somewhere on positive
sampling mass.  In Lean this is recorded generically as
\[
  \texttt{BitEncodedClassifierFamily.not\_badCodeAgreementBound\_of\_agrees\_on\_support}
\]
and, contrapositively, as
\[
  \texttt{BitEncodedClassifierFamily.exists\_pos\_mass\_disagreement\_of\_badCode\_agreementBound\_lt\_one}.
\]
The atomic special case is also pinned: for \(\mu=\delta_x\), any bad code that
agrees with the target at \(x\) has agreement mass exactly \(1\).  This is
recorded as
\[
  \texttt{BitEncodedClassifierFamily.not\_badCodeAgreementBound\_pure\_of\_agrees}
\]
and at the exact ZAB interfaces as
\[
  \texttt{not\_ExactZABERMRecoveryData\_of\_pure\_badCode\_agrees},
  \quad
  \texttt{not\_CanonicalZABERMRecoveryData\_of\_pure\_badCode\_agrees}.
\]
Equivalently, under a pure-atom bound with \(q<1\), any code that matches the
target at that atom must be the target function globally.  Thus finite-image
data for the ERM wrapper is still insufficient: the manuscript must also prove
a real support/disagreement theorem for the actual switched bad codes and the
measure being used.

\section{Exact Post-Switch Summary Compression}

The next optimistic cleanup is to isolate the exact object that still needs to
be compressed.  The module \texttt{ExactSwitchedFamily.lean} now makes that
object explicit.

The manuscript-visible post-switch datum is the exact local surface
\[
  u = (z,a,b).
\]
The Lean file separates:
\begin{enumerate}[leftmargin=1.5em]
  \item the full exact visible surface \(u=(z,a,b)\),
  \item the invariant projection \((z,a)\),
  \item the fork projection \(((z,a),b)\).
\end{enumerate}

The main positive result is not yet a proof that the actual switched family is
small.  It is the next honest theorem underneath that goal.

\begin{theorem}[Finite summary implies explicit bit budget]
Let \(G\) be an indexed Boolean predictor family on some input type \(U\).  If
\(G\) factors through a map
\[
  \mathrm{view} : U \to V
\]
to a finite summary surface \(V\), then \(G\) admits a concrete truth-table bit
budget of size \(|V|\).
\end{theorem}

\noindent
This is formalized as
\texttt{IndexedPredictorFamily.hasBitBudget\_card\_of\_factorsThrough\_fintype}.

So the optimistic burden is now cleaner.  It is no longer
\begin{quote}
compress the switched family somehow.
\end{quote}
It is:
\begin{quote}
prove that the switched family depends only on a genuinely small finite summary
surface.
\end{quote}

The exact post-switch specializations are now also explicit.

\begin{theorem}[Invariant-view compression kernel]
If the exact switched family on \(u=(z,a,b)\) factors through the invariant
projection \((z,a)\), then it has an exact visible compression target with bit
budget
\[
  |Z| \cdot 2^k.
\]
\end{theorem}

\noindent
This is formalized by
\texttt{card\_invariantVisiblePostSwitchSurface} together with
\texttt{exactVisibleCompressionTarget\_of\_factorsThrough\_invariant\_fintype}.

\begin{theorem}[Fork-view compression kernel]
If the exact switched family factors through the fork projection
\(((z,a),b)\), then it has an exact visible compression target with bit budget
\[
  |Z| \cdot 2^k \cdot 2^k.
\]
\end{theorem}

\noindent
This is formalized by
\texttt{card\_forkVisiblePostSwitchSurface} together with
\texttt{exactVisibleCompressionTarget\_of\_factorsThrough\_fork\_fintype}.

This does \emph{not} defeat the crux.  Those summary surfaces may still be far
too large.  But it does give the ground-up reduction we wanted: a successful
repair must now prove an actually small finite summary theorem, not merely a
slogan about switched locality.

The new module \texttt{GlobalWeaknessObstruction.lean} makes the first exact
Meredith-side pressure point equally explicit.  In the current Lean bridge,
\texttt{distinctionWeakness ev} and \texttt{nonDistinctionWeakness ev} are
global scalars attached to the whole finite evidence assignment \(ev\).  They do
not vary with the underlying state \(u \in U\).

\begin{theorem}[Weakness-only factoring is constant]
If a boolean predictor \(h : U \to \{\mathrm{false},\mathrm{true}\}\) factors
through either of the current global summaries
\[
  \mathrm{distinctionWeakness}(ev)
  \qquad \text{or} \qquad
  \mathrm{nonDistinctionWeakness}(ev),
\]
then \(h\) is constant on \(U\):
\[
  \forall u,v \in U,\; h(u) = h(v).
\]
\end{theorem}

\noindent
This is formalized as
\texttt{factorsThroughDistinctionWeakness\_constant} and
\texttt{factorsThroughNonDistinctionWeakness\_constant}.  So the current
weakness bridge is real semantic infrastructure, but by itself it does not yet
instantiate the same-route learning interface.  Any GSLT/weakness rescue that
still wants to count as the current switching/ERM route must add genuinely
per-instance observables or side information beyond these present global
summaries.

\section{Hyperseed as a Local-View Surface}

The existing Hyperseed formalization already provides a clean observer-relative
semantics.

Given a perspective \(P\), a budget \(B\), and a guard region \(G\), it defines
an \emph{available region}
\[
  \mathrm{availableRegion}(P,B,G)
  =
  \mathrm{observableUniverse}(P)
  \cap
  \mathrm{nearEurycosm}(P,B)
  \cap
  G.
\]

For the grassroots PNP program, this matters because it gives a principled way to
talk about what a decoder can even access.  Instead of treating locality as an
informal phrase, we can work with a typed input surface:
\[
  \mathrm{AvailableInput}(P,B,G)
  := \{x : \mathrm{World} \mid x \in \mathrm{availableRegion}(P,B,G)\}.
\]

The new Lean module proves the corresponding encoded-family bound on such an
available input surface:

\begin{theorem}[Available-region code-budget bound]
An \(s\)-bit classifier family over one Hyperseed available region still realizes
at most \(2^s\) classifiers.
\end{theorem}

\noindent
This is formalized as
\texttt{card\_realized\_availableFamily\_le\_of\_bitBudget}.

The theorem is intentionally modest.  It does not claim that the current local
view \emph{is} a Hyperseed perspective.  It says that Hyperseed already gives us a
good general language for bounded observational access, and that this language can
interoperate with the encoded-hypothesis-class infrastructure.

\section{Operational Distinction and Weakness}

My own mathematical and computationalist intuition is that if there is an elegant
route through the real blockers, it will likely come from a semantics of
\emph{distinguishability} rather than from ad hoc local-rule counting alone.
That is the reason the Meredith lane is worth treating as part of the grassroots
background theory, at least optimistically.

At the abstract level, the bridge has three layers:
\begin{enumerate}[leftmargin=1.5em]
  \item a GSLT supplies an operational state space together with its bisimulation
  quotient;
  \item context-decorated HML supplies a language of experiments or observations;
  \item quantale weakness measures the mass of distinction events on that quotient.
\end{enumerate}

In slogan form:
\begin{quote}
operational behavior gives the ontology, HML gives the observables, and weakness
turns distinguishability into evidence.
\end{quote}

This is now more than a philosophical analogy in the local Lean state.

\begin{theorem}[Depth-bounded observational equivalence]
For each modal depth \(n\), the Lean development defines a relation
\[
  \mathrm{hmlEquivUpTo}_n(t,u),
\]
meaning that no context-decorated HML formula of modal depth at most \(n\)
distinguishes \(t\) from \(u\).
\end{theorem}

\noindent
This is formalized in \texttt{GSLT/Logic/LogicalMetric.lean}.  The same file also
defines a binary-valued depth-\(n\) logical-distance approximation
\[
  d_n(t,u) \in \{0,1\},
\]
and proves the corresponding pseudoultrametric inequality.

\begin{theorem}[Finite-depth witness forces quotient-level distinction under adequacy]
Assume the sound direction of Meredith's adequacy theorem for a GSLT \(S\).  If
some depth-\(n\) context-decorated HML formula distinguishes \(t\) from \(u\),
then \(t\) and \(u\) are genuinely distinct in the bisimulation quotient of
\(S\).  Equivalently, if
\[
  d_n(t,u) = 1,
\]
then the bisimulation classes \([t]\) and \([u]\) are different.
\end{theorem}

\noindent
This is now formalized in \texttt{GSLT/Meredith/WeaknessBridge.lean} as
\texttt{distinguishingWitness\_implies\_distinguished\_of\_adequacySound} and
\texttt{logicalDistanceApprox\_eq\_one\_implies\_classes\_ne\_of\_adequacySound}.

\begin{theorem}[Weakness on the bisimulation quotient]
For a finite bisimulation quotient \(U\) equipped with a quantale-valued weight
assignment \(\mu\), the Lean bridge defines the weakness of the distinction event
\[
  \mathrm{weakness}_{\mu}(\mathrm{distinctionEvent}(U)),
\]
interpreting operational distinguishability as an evidence quantity.
\end{theorem}

\noindent
This is the content of \texttt{GSLT/Meredith/WeaknessBridge.lean}, building on
\texttt{GSLT/Meredith/Bisimulation.lean}.  In the finite probabilistic case, this
has the expected logical-entropy shape.

The new concrete step is that the abstract minimal-context interface now has an
actual process-calculus instance.

\begin{theorem}[Concrete rho-calculus minimal-context bridge]
The rho-calculus evaluation contexts from
\texttt{Languages/ProcessCalculi/RhoCalculus/Context.lean} now induce a concrete
\texttt{HasMinimalContexts} instance for the rho GSLT, and for each rho context
\(K\),
\[
  \mathrm{contextStep}(p,K,q)
  \iff
  p \mathrel{\rightsquigarrow[K]} q.
\]
\end{theorem}

\noindent
This is formalized in \texttt{GSLT/Meredith/RhoMinimalContext.lean}.  That file
also proves concrete HML diamonds for matching COMM interactions.

For the PNP grassroots route, the optimistic use of this material is not that it
proves the switching theorem by itself.  Rather, it offers a clean semantics for
coarse observational views.  If a decoder is only sensitive to finite-depth local
experiments, then one can hope to represent that by an HML fragment and then ask
how much evidence survives under the induced quotient.  That is exactly the sort
of elegant abstraction that may help disentangle the real compression/locality
story from the current proof's more brittle prose.  The bridge is now one step
tighter than before: finite-depth observational separation no longer just lives
as a metric-like approximation, but already certifies distinct quotient classes
once adequacy-soundness is available.

The honest limit is equally important: this is still background theory.  It does
\emph{not} yet produce the switched small class, the generalization theorem, or
the \(K_{\mathrm{poly}}\) contradiction step.  But it is now a real machine-checked
background layer, not merely an idea.
The new obstruction above sharpens this limit: the currently formalized
weakness summaries are quotient-level evidence scalars, not a per-instance
predictor surface.  So any semantic rescue that remains close to the
current switching/ERM route still needs an extra layer that turns operational
semantics into state-indexed encoded observables.

The next obstruction is sharper still.  The new module
\texttt{InfinitaryHMLObstruction.lean} proves:
\begin{enumerate}[leftmargin=1.5em]
  \item if \(\mathrm{MinimalContext}(S)\) is infinite, then \(\mathrm{HMLFormula}(S)\)
  is infinite;
  \item under the same hypothesis, even the depth-1 fragment
  \(\mathrm{DepthFragment}_1(S)\) is infinite;
  \item in the concrete rho instance, there is an explicit infinite ladder of
  certified minimal contexts, so the rho depth-1 HML fragment is already
  infinite;
  \item therefore no finite code space can surject onto the current exact
  depth-1 probe family, abstractly or in the concrete rho instance.
\end{enumerate}

So the current Meredith semantic layer is now pinned down even more tightly:
before one can talk about \(\mathrm{polylog}(m)\)-bit encoded predictor families,
one first needs a separate theorem that compresses or finitizes the relevant
observational probe family.  The present HML/context layer does not supply that
for free.

The Meredith-side trilemma is now simple.
\begin{enumerate}[leftmargin=1.5em]
  \item Use the current weakness summaries: then the predictor is constant.
  \item Use the current exact HML/context probes: then the observable family is
  already too large, even at depth \(1\).
  \item Use the current exact present-moment surface directly: then one already
  encounters powerset-like finite probe slices.
\end{enumerate}

So any close-form repair now has to do something genuinely new.  It is no
longer enough to say ``use current weakness'' or ``use current HML probes a bit
more cleverly.''  One needs a new compressed observational layer together with
an explicit theorem showing that the switched predictor factors through it.

The new module \texttt{PresentMomentShattering.lean} tightens this further on a
very concrete slice of the current rho semantics.  Let state \(i\) be a single
output on channel \(i\), and let environment \(S\) be a flat bag of input
guards on the channels in a finite list \(S\).  Then the file proves:
\begin{enumerate}[leftmargin=1.5em]
  \item channel \(i\) lies in
  \(\mathrm{surfaceChannels}(\mathrm{state}_i,\mathrm{env}_S)\) iff \(i \in S\);
  \item the concrete external-present-moment pair
  \((\mathrm{par}(\mathrm{env}_S,\Box), i)\) is present iff \(i \in S\).
\end{enumerate}

So even before passing to the exact HML layer, the currently formalized
state-indexed observational surface already realizes arbitrary finite
membership patterns on a natural probe family.  This is not yet a full
``depth-1 HML shattering'' theorem, but it means the large-observable-surface
pressure is already visible one layer earlier in the exact present-moment
semantics.  In fact, the file now packages this as an injective finite-signature
map \(S \mapsto (i \mapsto \mathrm{observable}_S(i))\) on
\(\mathrm{Finset}(\mathrm{Fin}\,n)\): distinct finite probe sets yield distinct
exact present-moment signatures on the \(n\) probe states.  Any close-form
Meredith rescue now has to compress not only an infinitary exact HML layer, but
already a powerset-like concrete probe family in the rho present-moment surface.

\section{Kolmogorov Foundations Already Present}

The local repository also already contains a substantial Kolmogorov-complexity seed
under \texttt{Mettapedia/Computability/KolmogorovComplexity/}.

In particular, \texttt{Basic.lean} defines:
\begin{enumerate}[leftmargin=1.5em]
  \item binary strings as \texttt{List Bool},
  \item partial algorithms,
  \item plain Kolmogorov complexity \(C[U](x)\),
  \item a universal algorithm built from mathlib's partial-recursive code,
  \item the usual invariance theorem for universal algorithms.
\end{enumerate}

This is not yet the time-bounded conditional complexity \(K_{\mathrm{poly}}\)
that the present route uses.  But it matters for the grassroots route because it
means the eventual \(K_{\mathrm{poly}}\) development should be an extension of an
existing computability layer rather than a fresh bespoke definition dropped into a
late-stage contradiction proof.

In practical terms, the next computability milestone is:

\begin{quote}
define conditional and polynomial-time-bounded complexity on top of the existing
universal-machine layer, instead of axiomatizing them only at the paper level.
\end{quote}

\section{What the Grassroots Program Needs Next}

The current background theory suggests a fairly crisp optimistic roadmap.

\subsection{Finite-Class Learning Layer}

The ERM existence layer, finite uniform consistency bound, threshold recovery
corollary, rational uniform-rate layer, miss-ratio success bound, and weighted
one-predictor and finite-family weighted bounds are now in place.  The next
finite-class learning
steps are:
\begin{enumerate}[leftmargin=1.5em]
  \item a generalization bound stated in terms of \(\log |H|\),
  \item a probability interface for i.i.d. or distributional assumptions at the family level,
  \item concentration-style control beyond the exact finite-space formulas,
  \item a clean interface separating statistical assumptions from coding facts.
\end{enumerate}

In other words, the current Lean state already handles:
\begin{enumerate}[leftmargin=1.5em]
  \item empirical error on finite samples,
  \item existence of ERM predictors in nonempty finite encoded families,
  \item exact-fit transfer when some code in the family is sample-consistent,
  \item a finite uniform bound on the deceptive sample set,
  \item exact ERM recovery of the target off that deceptive set,
  \item a strict-cardinality threshold guaranteeing the existence of a good sample,
  \item a rational upper bound on the deceptive sample rate,
  \item a rational lower bound on the non-deceptive sample rate,
  \item a rational lower bound on the exact recovery rate,
  \item strict positivity of that exact recovery rate under the threshold,
  \item explicit miss-ratio lower bounds of the form \(1 - |H|\rho_X^m\),
  \item an exact weighted one-predictor law of the form \(\mathrm{mass} = \alpha^m\),
  \item a weighted family union bound of the form
  \(\mathrm{deceptiveMass} \le \sum_{h \in H_{\mathrm{bad}}}\alpha_h^m\),
  \item a weighted exact recovery bound of the form
  \(\mathrm{recoveryMass} \ge 1 - \sum_{h \in H_{\mathrm{bad}}}\alpha_h^m\).
\end{enumerate}

\noindent
What remains is to lift that finite uniform rate layer into a real finite-class
generalization theorem.

\subsection{Concrete Decoder Subclasses}

The encoded-family theorem is only an interface.  It still needs a nontrivial
source of small classes.  The likely optimistic candidates are concrete local
decoders on tree-like neighborhoods:
\begin{enumerate}[leftmargin=1.5em]
  \item belief-propagation style update rules,
  \item GF(2)-linear constraint summaries from the VV layer,
  \item bounded-depth compositions of those primitives.
\end{enumerate}

The first concrete instance of that second lane now exists.  The new module
\texttt{AffineColumnFamily.lean} formalizes affine GF(2)-style predictors on the
retained VV column bits and proves:
\begin{quote}
if the exact switched family collapses to affine parity tests on the column
view \(a\), then the exact visible compression target is only \(k+1\) bits.
\end{quote}

\noindent
The key formal statement is
\texttt{exactVisibleCompressionTarget\_of\_realizedByExactAffineColumnFamily}.
This does not prove that the real switched family is affine.  It does give the
optimistic route its first genuinely nontrivial concrete landing zone below the
full truth-table class.

The next concrete step is now also formalized.  The new module
\texttt{AffineFeatureFamily.lean} allows a bounded number \(r\) of affine
GF(2)-style VV-column features, followed by an arbitrary truth table on the
resulting \(r\) feature bits.  Its exact raw bit budget is
\[
  r(k+1) + 2^r.
\]
The key formal statement is
\texttt{exactVisibleCompressionTarget\_of\_realizedByExactAffineFeatureFamily}.
So if the switched family collapses not to a single affine probe but to a
bounded affine-feature algebra, the code budget is still explicit and still far
below the full truth-table class on \(u=(z,a,b)\).

The next concrete step is sharper and more useful.  The new module
\texttt{SparseThresholdAffineFamily.lean} keeps the same \(r\) affine
GF(2)-style VV-column features but replaces the full truth table by a much
smaller combiner: an \(r\)-bit mask selecting active features and an \(r\)-bit
threshold code.  Its exact raw bit budget is now linear:
\[
  r(k+3).
\]
The key formal statement is
\texttt{exactVisibleCompressionTarget\_of\_realizedByExactSparseThresholdAffineFamily}.
So if the switched family collapses to a sparse-threshold affine algebra on the
retained VV column view, the optimistic route now has a concrete candidate
class whose coding cost no longer contains the exponential \(2^r\) blow-up.

If one of these affine / sparse-threshold families can be shown to encode the
actual switched predictors uniformly in \(\mathrm{polylog}(m)\) bits, then the
switching theorem has a plausible landing zone.

The next refinement is now also formalized.  The new module
\texttt{AffineDecisionListFamily.lean} keeps the same \(r\) affine GF(2)-style
VV-column features, but replaces the combiner by a fixed-order decision list:
each feature carries one response bit, and there is one default output bit if
no feature fires.  Its exact raw bit budget is
\[
  r(k+2) + 1.
\]
The key formal statement is
\texttt{exactVisibleCompressionTarget\_of\_realizedByExactAffineDecisionListFamily}.
So if the switched family collapses to an ordered affine decision-list algebra
on the retained VV column view, the optimistic route now has an even smaller
explicit candidate class than the sparse-threshold family.

The next layer is now also formalized.  The new module
\texttt{ExactAffineRecovery.lean} turns those four exact-surface candidate
classes into direct weighted finite-sampling recovery theorems.  Once the
target lies in one of the concrete exact-surface affine families, the existing
recovery backbone from \texttt{SameRouteInterface.lean} applies immediately
with one of the four explicit budgets
\[
  k+1, \qquad r(k+1)+2^r, \qquad r(k+3), \qquad r(k+2)+1.
\]
So the current small-class ladder is no longer only a compression interface.
It is now a direct recovery interface too.

The next big refinement is now also formalized.  The new module
\texttt{SharedAffineFeatureFamilies.lean} isolates the regime where one fixed
affine feature basis is shared across the switched family and only the
downstream combiner varies.  In that case the feature extractor no longer has
to be re-encoded per predictor, so the exact visible compression target drops
to the combiner alone:
\[
  2^r, \qquad 2r, \qquad r+1
\]
for the shared truth-table, shared sparse-threshold, and shared decision-list
regimes respectively.  The same module also proves the corresponding weighted
exact-recovery lower bounds on the exact post-switch surface.

This is the first genuinely sharp improvement in the optimistic lane.  If the
real switched family uses a common low-complexity feature extractor and varies
only in a light combiner, then the coding burden is much smaller than in the
previous per-predictor affine-family models.

The next concrete step is now also formalized.  The new module
\texttt{ExactABDecisionListFamily.lean} fixes the shared extractor to the raw
exact visible bits themselves: concatenate the retained VV column bits \(a\)
with the side-channel bits \(b\), then run a fixed-order decision list on those
\(2k\) raw bits.  This yields a completely concrete exact-surface candidate
class with budget
\[
  2k + 1.
\]
The key point is that the feature extractor is no longer existential or
parameterized.  It is the literal raw exact visible surface restricted to the
bit coordinates \((a,b)\).  The same module also proves the corresponding
weighted exact-recovery lower bound for that family.

The route is now factored one step more cleanly.  The new modules
\texttt{ABVisibleSurface.lean} and \texttt{ABDecisionListRoute.lean} isolate the
reduced raw visible surface \((a,b)\) as an explicit quotient candidate and
prove the pullback statement one actually wants: if the switched family factors
through that reduced surface and the reduced family is realized by fixed-order
decision lists there, then the exact-surface compression and weighted recovery
theorems follow immediately.  This is also the natural insertion point for any
future observational-equivalence argument: a semantic quotient can help justify
the map to \((a,b)\), but it does not need to alter the learning-theoretic
backbone once that map is fixed.  The same route file now also proves the
slightly stronger lift: if an exact-surface family is invariant under the
\((a,b)\) quotient, then a canonical reduced family can be reconstructed and
the same \(2k+1\)-bit decision-list target follows from that invariance alone.

The matching strict-subclass fact is now machine-checked as well.
\texttt{ABDecisionListObstruction.lean} defines the fixed-order decision-list
bit family directly on the reduced raw surface and proves that realization by
that family gives a \(2k+1\)-bit budget on \(\texttt{ABVisibleSurface}\ k\).
For \(k>0\), Lean proves
\[
  2k+1 < 2^{2k} = |\texttt{ABVisibleSurface}\ k|.
\]
Hence the raw decision-list route cannot realize the full Boolean rule family
on \((a,b)\).  The same obstruction lifts through the exact-surface pullback:
if the exact switched family factors through \((a,b)\) but the reduced family is
still fully expressive, then the raw exact-surface decision-list class cannot
realize it.  This is grassroots progress rather than a final blocker: it says
exactly which finite candidate class has been bounded, and leaves broader
randomized, approximate, and \(Kpoly\) bridges outside the claim.

The raw \((a,b)\) route is now richer than a single fixed-order decision list.
The new module \texttt{ExactABAffineFamilies.lean} reuses the existing affine
machinery on the same raw extractor and proves exact-surface compression and
weighted recovery bounds for three further candidate classes on the raw
\((a,b)\) bits:
\[
  r((2k)+1)+2^r, \qquad r((2k)+3), \qquad r((2k)+2)+1
\]
for bounded affine features, sparse-threshold affine predictors, and affine
decision lists respectively.  So the reduced raw surface is now a whole
candidate-family ladder, not just one isolated family.

The next refinement is now in place as well.  The new module
\texttt{SharedExactABFeatureFamilies.lean} adds the shared-basis collapse on
that same raw \((a,b)\) surface.  If one fixed affine basis on the \(2k\) raw
bits is reused across the whole switched family and only the downstream
combiner varies, the exact visible compression target drops all the way to
combiner-only budgets:
\[
  2^r, \qquad 2r, \qquad r+1
\]
for arbitrary truth tables, sparse-threshold combiners, and fixed-order
decision lists respectively.  So the raw \((a,b)\) quotient now supports both a
per-predictor affine ladder and a sharper shared-basis ladder.

That bridge is now factored explicitly too.  The new module
\texttt{SharedABFeatureRoutes.lean} proves that if the switched family factors
through the reduced raw surface \((a,b)\) and the reduced family is realized by
one of those shared-basis raw affine classes, then the exact-surface family
inherits the same combiner-only budget immediately.  This is the sharpest
current handoff point between any future observational quotient theorem and the
small-class recovery machinery.  The same file now also proves the invariance
version: if the exact-surface family is already invariant under the raw
\((a,b)\) quotient, then the canonical reduced family inherits those same
shared-basis combiner-only budgets without any separately supplied
factorization witness.

The route-coverage map now records the matching strict-subclass obstruction as
well.  The module \texttt{SharedABFeatureObstruction.lean} proves that if the
reduced raw family on \((a,b)\) is still surjective onto all Boolean rules,
then no fixed shared raw affine basis can realize it below the reduced
truth-table threshold.  Lean proves this for the three combiner budgets
\[
  2^r, \qquad 2r, \qquad r+1
\]
corresponding to arbitrary affine-feature truth tables, sparse-threshold
combiners, and fixed-order decision lists.  The same result is lifted back to
the exact surface: if the exact switched family factors through \((a,b)\) and
the reduced family is fully expressive, then the corresponding shared exact
class cannot realize it below those budgets.  This adds honest narrow entries
to the \texttt{PNPKpolySubrepairClass} ledger alongside the finite-image
budget, representative-index, quotient-code, transport, and reduced-AB
finite-image/quotient pullback entries, while leaving the broad
\texttt{kpolyCompressionBridge} open; it
rules out full reduced-rule expressivity for named shared classes, not every
possible manuscript-faithful small-class theorem.

The target statement itself is now packaged in that style too.  The new module
\texttt{SharedABTargetInterface.lean} bundles the current route bottleneck into
explicit data packages: one shared raw \((a,b)\) feature basis, one quotient
invariance hypothesis, and one realization hypothesis for the lifted reduced
family.  From that package it derives the exact visible compression target with
budget \(2^r\), \(2r\), or \(r+1\) depending on the combiner class.  This is
the cleanest current formal statement of what the optimistic route actually
needs.

There is now also a clean backward step from that generic target statement back
toward the most concrete current route.  The module
\texttt{CanonicalABTargetRoute.lean} specializes the shared raw \((a,b)\)
feature basis to the canonical coordinate probes on the raw visible bits.  In
that specialization, the generic shared-basis decision-list target collapses
exactly to the older raw \((a,b)\) decision-list route.  So the current formal
burden can now be read in two equivalent ways:
\begin{enumerate}[leftmargin=1.5em]
  \item as a generic shared-basis target interface on the reduced raw surface,
  \item or as quotient invariance plus realization by fixed-order decision
  lists on the literal raw visible bits.
\end{enumerate}
This is the clearest backward-chaining reach from the cleaned-up target
statement back to the existing concrete work.

That concrete work is now packaged directly as well.  The new module
\texttt{CanonicalABCandidateInterface.lean} takes the canonical raw-bit route
and bundles it into one explicit candidate-data object: quotient invariance
under the raw \((a,b)\) projection together with realization of the lifted
reduced family by fixed-order decision lists on the raw visible bits.  From
that single package it derives both the exact visible compression target with
budget \(2k+1\) and, for each indexed predictor in the family, the matching
weighted exact-recovery lower bound through the already-existing raw decision
list bit family.  So the optimistic burden is now even sharper: produce one
such candidate-data witness, then prove the needed agreement-mass estimate.

That burden has now been reduced one step further.  The new module
\texttt{CanonicalABCodeWitness.lean} proves that if the switched family is
exactly the family induced by an explicit assignment of one raw \((a,b)\)
decision-list code to each index, then the canonical candidate-data package,
the \(2k+1\) compression target, and the per-index recovery theorem all follow
automatically.  So the route bottleneck is now close to its simplest concrete
form: show that the real switched family equals one such code-induced family,
then prove the corresponding agreement-mass estimate.

The last generic layer is now also in place.  The new module
\texttt{CanonicalABRecoveryInterface.lean} packages exactly the remaining
generic information needed by the route:
\begin{enumerate}[leftmargin=1.5em]
  \item one explicit raw \((a,b)\) decision-list code assignment per index,
  \item equality between the switched family and the induced code family,
  \item one agreement-mass upper bound.
\end{enumerate}
From that single recovery-data package, the \(2k+1\) compression target and the
weighted exact-recovery lower bound both follow automatically.  At that point,
the remaining work is no longer background theory; it is the actual
candidate-specific switched-family theorem.

The next manuscript-facing extension is now in place as well.  The paper's real
local input is not just \((a,b)\), but \(u=(z,a,b)\), where \(z\) comes from one
fixed shared extractor on the masked block.  The new module
\texttt{ExactZABDecisionListFamily.lean} formalizes the corresponding candidate
class directly: for any shared extractor \(\mathsf{zfeat}: Z \to \{0,1\}^r\), one
concatenates \(\mathsf{zfeat}(z)\), \(a\), and \(b\), then applies the same
fixed-order decision-list family as before.  The exact-surface budget is then
\(r+2k+1\), and the matching weighted exact-recovery bound is proved on that
larger visible surface.

That new route is then packaged in \texttt{ExactZABTargetInterface.lean}.  The
burden is exactly what one would want it to be: prove that the switched family
is realized by fixed-order decision lists on \((\mathsf{zfeat}(z),a,b)\), then
prove one bad-code agreement-mass estimate.  No extra generic machinery is
needed beyond that package.  The same package now also exposes the finite-image
content of the route explicitly: from the target data it derives a finite
predictor cover, finite representative-index cover, and finite quotient-code
presentation of size at most \(2^{r+2k+1}\).  Thus the concrete exact-ZAB route
does not hide the current \(Kpoly\) repair target inside the compression-target
abbreviation.

Finally, the manuscript's constructive wrapper now plugs into the route
directly.  The new module \texttt{BitFamilyERM.lean} turns any fixed bit-encoded
local class into an indexed ERM-selected family together with its selected code
at each sample.  Then \texttt{ExactZABERMRoute.lean} specializes that statement
to the shared \((\mathsf{zfeat}(z),a,b)\) decision-list class: if the wrapper is
honestly specified as ERM over that class, then the selected code itself is the
route witness, and the compression target follows immediately.  The same ERM
route now exposes the finite predictor cover, representative-index cover, and
quotient-code presentation for the ERM-selected family itself.  At that point,
the remaining burden is even sharper: identify the right shared extractor
\(\mathsf{zfeat}\), show the wrapper is using that class, and control the
agreement mass of the bad codes.

There is now one more collapse on top of that manuscript-facing route.  The new
module \texttt{SharedExactZABFeatureFamilies.lean} assumes that the local bits
\((\mathsf{zfeat}(z),a,b)\) first pass through one fixed shared affine-feature
basis of width \(p\), and only the final combiner varies across predictors.  On
that shared-basis route, the exact-surface budgets fall to \(2^p\) for arbitrary
truth tables, \(2p\) for sparse thresholds, and \(p+1\) for fixed-order
decision lists.  Then \texttt{SharedExactZABTargetInterface.lean} packages the
corresponding target and recovery interfaces, and
\texttt{SharedExactZABERMRoute.lean} shows that if the wrapper is honestly ERM
over that shared-basis decision-list class, then the selected code itself again
is the route witness, now with the sharper \(p+1\) budget.  The shared-basis
exact route now also exposes the same finite-image data as the direct exact
route.  \texttt{SharedExactZABTargetInterface.lean} derives finite predictor
covers, representative-index covers, and quotient-code presentations for the
affine-feature, sparse-threshold, and shared decision-list variants.  Then
\texttt{SharedExactZABFeatureERMInterface.lean} packages the final ERM-facing
affine-feature and sparse-threshold forms, and
\texttt{SharedExactZABERMInterface.lean} packages the final ERM-facing shared
decision-list form.  So at the shared-basis endpoint, one sample assignment,
one ERM-equality witness, and one bad-code agreement bound already yield both
the weighted recovery inequality and the explicit finite-image payload.

So the remaining burden is now split cleanly into two candidate routes on the
real local surface \(u=(z,a,b)\):
\begin{enumerate}[leftmargin=1.5em]
  \item the direct decision-list route on \((\mathsf{zfeat}(z),a,b)\), with budget \(r+2k+1\),
  \item the shared-basis route, where one fixed affine basis on \((\mathsf{zfeat}(z),a,b)\) reduces the decision-list budget to \(p+1\).
\end{enumerate}
The second route is now the sharper optimistic landing point.

The direct route has now been backward-chained as far as it can go generically.
The new module \texttt{CanonicalZABTargetRoute.lean} specializes the
shared-basis statement to the canonical coordinate basis on the raw exact bits,
so the shared-basis route collapses back to the direct exact decision-list
route.  Then \texttt{CanonicalZABCandidateInterface.lean} packages that direct
route as one concrete candidate object, \texttt{CanonicalZABCodeWitness.lean}
reduces it to one explicit exact decision-list code assignment per index, and
\texttt{CanonicalZABRecoveryInterface.lean} packages the last direct-route
ingredient: one agreement-mass upper bound.  So even on the less sharp
\(r+2k+1\) route, the remaining burden is now completely explicit.  Moreover,
\texttt{CanonicalZABERMRoute.lean} removes the last administrative gap between
the honest ERM wrapper and that final direct route: once the wrapper is really
ERM over the raw exact \((\mathsf{zfeat}(z),a,b)\) decision-list class, the
canonical code witness is exactly the ERM-selected code.  The new module
\texttt{CanonicalZABERMInterface.lean} packages the same conclusion in the most
manuscript-facing form: one sample assignment, one proof that the switched
family is exactly that ERM wrapper, and one bad-code agreement bound.  The
canonical candidate, explicit-code, recovery, and final ERM layers now carry
the same finite predictor cover, representative-index cover, and quotient-code
objects directly.  This discharges the finite-image bookkeeping for the actual
canonical exact decision-list family conditional on an explicit code assignment
or equality to the ERM wrapper; it does not prove that the manuscript's broader
switched predictor family has such a realization.

There is now also one fully concrete specialization of that exact route.
\texttt{BitVecZABVisibleSurface.lean} fixes the case where the local datum is
already a bitvector \(z : \mathrm{BitVec}\, r\), so the full visible input
\((z,a,b)\) becomes one raw bitvector of length \(r+2k\).  Then
\texttt{BitVecZABERMInterface.lean} specializes the final exact ERM interface
to the identity extractor on that raw full visible bit surface.  So in that
case the remaining burden is no longer “choose \(\mathsf{zfeat}\)”; it is only
“choose the samples, prove equality to the ERM wrapper, and control the bad
codes.”  The new module \texttt{ProjectedZABERMInterface.lean} narrows the same
burden in a different concrete way: when \(z\) is bit-valued, it is enough to
choose finitely many coordinates of \(z\) as the local summary extractor.
Both concrete ERM specializations now also expose the three finite-image
objects directly, so the remaining manuscript obligation is not hidden behind
the identity or projection wrappers.

\subsection{Bridging to the Full PNP Route}

Only after those layers are built does it make sense to reconnect to the full
proof program:
\begin{enumerate}[leftmargin=1.5em]
  \item define the distributional ensemble cleanly,
  \item define local views as typed accessible surfaces,
  \item prove the switched family lies in a small encoded subclass,
  \item connect success bounds to a properly formalized \(K_{\mathrm{poly}}\),
  \item finally compare with the upper bound under \(P=NP\).
\end{enumerate}

\section{Current Status}

At present, the grassroots side has several real machine-checked gains and one
major pre-existing base layer:

\begin{center}
\begin{tabular}{lll}
\toprule
\textbf{Layer} & \textbf{Status} & \textbf{Artifact} \\
\midrule
encoded hypothesis families & proved & \texttt{HypothesisClass.lean} \\
finite-class ERM kernel & proved & \texttt{FiniteERM.lean} \\
finite uniform consistency bound & proved & \texttt{FiniteConsistencyBound.lean} \\
finite uniform recovery threshold & proved & \texttt{FiniteUniformRecovery.lean} \\
bit-family uniform recovery threshold & proved & \texttt{BitFamilyUniformRecovery.lean} \\
binary bit-family recovery threshold & proved & \texttt{BitFamilyUniformRecoveryBinary.lean} \\
finite uniform rate layer & proved & \texttt{FiniteUniformRate.lean} \\
finite uniform non-deceptive and recovery success bounds & proved & \texttt{FiniteUniformSuccess.lean} \\
Boolean quantitative bit-family non-deceptive, positive-margin, and recovery bounds & proved & \texttt{BitFamilyUniformSuccessBinary.lean} \\
clocked \(Kpoly\) construction, realization interface, and finite success-rate/good-sample bridge & proved & \texttt{ClockedKpolyBridge.lean} \\
clocked \(Kpoly\) gap assessment: payload equals finite predictor-image obligation & proved & \texttt{ClockedKpolyGapAssessment.lean} \\
clocked \(Kpoly\) exact ZAB ERM closure and bare-interface obstruction & proved conditional closure plus no-go & \texttt{ClockedKpolyActualGapClosure.lean} \\
actual switched local normal form: ZAB ERM closure and locality-only obstruction & proved minimal-assumption closure plus no-go & \texttt{ActualSwitchedLocalFamily.lean} \\
actual switched local strengthened structure: bounded-code and ZAB decision-list finite-image closure plus full-rule exclusions & proved minimal missing assumption & \texttt{ActualSwitchedLocalStructure.lean} \\
actual switched local ERM construction: concrete \(zOf,aOf,bOf,zfeat,samples\) wrapper & proved construction satisfies ZAB data, bit-code data, finite cover, and clocked payload & \texttt{ActualSwitchedLocalERMConstruction.lean} \\
actual switched local code construction: concrete \(zOf,aOf,bOf,zfeat,codes\) wrapper & proved construction satisfies ZAB data, bit-code data, finite cover, and clocked payload & \texttt{ActualSwitchedLocalCodeConstruction.lean} \\
actual switched local plug-in lookup endpoint & proved full-rule obstruction below truth-table budget & \texttt{ActualSwitchedLocalPluginLookupObstruction.lean} \\
actual switched local plug-in majority endpoint & proved full-rule obstruction below truth-table budget & \texttt{ActualSwitchedLocalPluginMajorityObstruction.lean} \\
actual switched local finite-sample plug-in majority endpoint & proved full-rule obstruction below truth-table budget & \texttt{ActualSwitchedLocalPluginSampleMajorityObstruction.lean} \\
actual switched local bounded-sample plug-in majority endpoint & proved exact surjectivity threshold at visible alphabet size & \texttt{ActualSwitchedLocalBoundedSampleMajorityObstruction.lean} \\
finite weighted one-predictor agreement law & proved & \texttt{FiniteIIDAgreement.lean} \\
finite weighted family union bound & proved & \texttt{FiniteIIDFamilyBound.lean} \\
finite weighted recovery bounds & proved & \texttt{FiniteIIDRecovery.lean} \\
same-route switching/ERM certificate & proved & \texttt{SameRouteInterface.lean} \\
same-route epsilon recovery certificate & proved & \texttt{BitFamilyWeightedRecovery.lean} \\
atomic bad-code agreement obstruction & proved & \texttt{BadCodeAgreementObstruction.lean} \\
exact post-switch finite-summary kernel & proved & \texttt{ExactSwitchedFamily.lean} \\
exact visible predictor-image budget equivalence & proved & \texttt{ExactVisibleImageBudget.lean} \\
reduced raw visible surface $(a,b)$ & proved & \texttt{ABVisibleSurface.lean} \\
affine GF(2)-style column candidate class & proved & \texttt{AffineColumnFamily.lean} \\
bounded affine-feature candidate class & proved & \texttt{AffineFeatureFamily.lean} \\
sparse-threshold affine candidate class & proved & \texttt{SparseThresholdAffineFamily.lean} \\
affine decision-list candidate class & proved & \texttt{AffineDecisionListFamily.lean} \\
exact-surface affine recovery ladder & proved & \texttt{ExactAffineRecovery.lean} \\
shared-basis affine family ladder & proved & \texttt{SharedAffineFeatureFamilies.lean} \\
raw exact-surface decision-list family on $(a,b)$ & proved & \texttt{ExactABDecisionListFamily.lean} \\
pullback route from reduced $(a,b)$ families & proved & \texttt{ABDecisionListRoute.lean} \\
raw $(a,b)$ decision-list strict-subclass obstruction & proved & \texttt{ABDecisionListObstruction.lean} \\
affine candidate ladder on raw $(a,b)$ bits & proved & \texttt{ExactABAffineFamilies.lean} \\
shared-basis affine ladder on raw $(a,b)$ bits & proved & \texttt{SharedExactABFeatureFamilies.lean} \\
quotient-route transfer to shared raw $(a,b)$ ladders & proved & \texttt{SharedABFeatureRoutes.lean} \\
shared raw $(a,b)$ strict-subclass and exact-pullback obstruction & proved & \texttt{SharedABFeatureObstruction.lean} \\
unified shared raw $(a,b)$ target interface & proved & \texttt{SharedABTargetInterface.lean} \\
canonical raw $(a,b)$ specialization of the target interface & proved & \texttt{CanonicalABTargetRoute.lean} \\
canonical raw $(a,b)$ candidate-data interface with recovery consequence & proved & \texttt{CanonicalABCandidateInterface.lean} \\
explicit canonical raw $(a,b)$ code-witness reduction & proved & \texttt{CanonicalABCodeWitness.lean} \\
final canonical raw $(a,b)$ recovery interface & proved & \texttt{CanonicalABRecoveryInterface.lean} \\
ERM transport layer for fixed bit-encoded local classes & proved & \texttt{BitFamilyERM.lean} \\
exact decision-list family on $(zfeat(z),a,b)$ & proved & \texttt{ExactZABDecisionListFamily.lean} \\
shared $(zfeat(z),a,b)$ target and recovery interfaces with finite-image presentations & proved & \texttt{ExactZABTargetInterface.lean} \\
ERM specialization of the shared $(zfeat(z),a,b)$ route with finite-image objects & proved & \texttt{ExactZABERMRoute.lean} \\
shared-basis affine ladder on $(zfeat(z),a,b)$ & proved & \texttt{SharedExactZABFeatureFamilies.lean} \\
shared-basis $(zfeat(z),a,b)$ target and recovery interfaces with finite-image presentations & proved & \texttt{SharedExactZABTargetInterface.lean} \\
final ERM interfaces for shared-basis affine-feature and sparse-threshold exact routes with finite-image objects & proved & \texttt{SharedExactZABFeatureERMInterface.lean} \\
ERM specialization of the shared-basis $(zfeat(z),a,b)$ route & proved & \texttt{SharedExactZABERMRoute.lean} \\
final ERM interface for the shared-basis exact $(zfeat(z),a,b)$ decision-list route with finite-image objects & proved & \texttt{SharedExactZABERMInterface.lean} \\
canonical specialization of the shared-basis exact route & proved & \texttt{CanonicalZABTargetRoute.lean} \\
canonical direct exact-route candidate interface with finite-image objects & proved & \texttt{CanonicalZABCandidateInterface.lean} \\
explicit direct exact-route code-witness reduction with finite-image objects & proved & \texttt{CanonicalZABCodeWitness.lean} \\
final direct exact-route recovery interface with finite-image objects & proved & \texttt{CanonicalZABRecoveryInterface.lean} \\
ERM-to-canonical bridge for the exact $(zfeat(z),a,b)$ route with finite-image objects & proved & \texttt{CanonicalZABERMRoute.lean} \\
final ERM interface for the canonical exact $(zfeat(z),a,b)$ route with finite-image objects & proved & \texttt{CanonicalZABERMInterface.lean} \\
raw full visible bit surface when $z : \mathrm{BitVec}\, r$ & proved & \texttt{BitVecZABVisibleSurface.lean} \\
identity-extractor ERM interface on the raw full visible bit surface with finite-image objects & proved & \texttt{BitVecZABERMInterface.lean} \\
coordinate-projection ERM interface on bit-valued local data with finite-image objects & proved & \texttt{ProjectedZABERMInterface.lean} \\
weakness-globality obstruction & proved & \texttt{GlobalWeaknessObstruction.lean} \\
infinitary HML obstruction & proved & \texttt{InfinitaryHMLObstruction.lean} \\
present-moment shattering slice & proved & \texttt{PresentMomentShattering.lean} \\
depth-bounded HML distinction kernel & proved & \texttt{LogicalMetric.lean} \\
GSLT bisimulation-to-weakness bridge & proved & \texttt{WeaknessBridge.lean} \\
concrete rho minimal-context instance & proved & \texttt{RhoMinimalContext.lean} \\
Hyperseed available-region reuse & proved & \texttt{HypothesisClass.lean} + \texttt{Hyperseed} \\
plain Kolmogorov infrastructure & already present & \texttt{KolmogorovComplexity/Basic.lean} \\
finite-class generalization bounds & not yet formalized & next target \\
\(K_{\mathrm{poly}}\) layer & not yet formalized & later target \\
\bottomrule
\end{tabular}
\end{center}

This is not the final proof.  It is the right sort of groundwork.

\section{Conclusion}

The grassroots program should stay separate from the crux program.

The crux note asks where the present paper overreaches.  This note asks what
general machinery we would still want even if the overreach is patched later.
On that constructive side, the current Lean result is that small encoded
hypothesis families, finite empirical-risk minimizers, finite deceptive-sample
bounds, finite uniform recovery thresholds, rational exact-recovery rates,
explicit miss-ratio success bounds, weighted one-predictor agreement laws,
weighted family deceptive-mass bounds, weighted exact recovery bounds,
same-route epsilon recovery certificates, exact post-switch finite-summary
compression kernels, one concrete affine GF(2)-style
candidate family, one bounded affine-feature candidate family, and
observer-relative input surfaces can already be handled cleanly inside one
coherent Mettapedia/mathlib formal background layer.  Optimistically, the same is
now becoming true for the observational-semantics side: finite-depth HML
distinction, weakness on operational quotients, and a concrete rho-calculus
minimal-context instance are landing as reusable infrastructure rather than as
one-off prose.  But the current Meredith bridge is now also honest about its own
shape: it produces global evidence summaries, not yet a nontrivial per-state
learning surface.  And even once one moves from those global summaries to the
HML/context layer, the current observable family is infinitary by default and
already concretely infinitary in the rho instance.

That is enough to sharpen the next optimistic task.  We should not try to prove
\(P \neq NP\) in one jump.  We should first prove that the intended switched
predictors live in a genuinely small encoded class.  The current best concrete
landing point is now sharper than the earlier raw \((a,b)\) route and sharper
than the first direct \((\mathsf{zfeat}(z),a,b)\) route: prove that the real
wrapper operates inside one shared-basis decision-list class on
\((\mathsf{zfeat}(z),a,b)\), and then prove the corresponding bad-code
agreement bound.  The direct \(r+2k+1\) route remains available, but the
shared-basis \(p+1\) route is now the best optimistic target.  Still, the
direct route is no longer vague: it has been reduced all the way to an explicit
code assignment plus an agreement bound.  Only after one of those routes is
instantiated should the rest of the argument be lifted on top of the
foundation.

\end{document}
